\documentclass[11pt,a4paper]{article}
\usepackage[T1]{fontenc}
\usepackage[utf8]{inputenc}
\usepackage{lmodern,amsmath,amssymb,amsthm,mathtools}
\usepackage[margin=26mm,headheight=15pt]{geometry}
\usepackage{microtype,booktabs,longtable,tabularx,array,enumitem}
\usepackage[dvipsnames]{xcolor}
\usepackage{listings,xurl,hyperref,fancyhdr}
\hypersetup{colorlinks=true,linkcolor=MidnightBlue,citecolor=MidnightBlue,urlcolor=MidnightBlue,pdftitle={Continuous mathematics in iPRA: measures, stochastic operators, and homological transfer},pdfauthor={Research prepared for Vladimir Reshetnikov}}
\newcommand{\N}{\mathbb N}
\newcommand{\Q}{\mathbb Q}
\newcommand{\R}{\mathbb R}
\newcommand{\C}{\mathbb C}
\newcommand{\ipra}{\mathsf{iPRA}}
\newcommand{\eps}{\varepsilon}
\newcommand{\norm}[1]{\left\lVert#1\right\rVert}
\newcommand{\abs}[1]{\left|#1\right|}
\newcommand{\E}{\mathbb E}
\newcommand{\id}{\operatorname{id}}
\newcommand{\clip}{\operatorname{clip}}
\newcommand{\dist}{\operatorname{dist}}
\newcommand{\supp}{\operatorname{supp}}
\newcommand{\PR}{\mathsf{PR}}
\newcommand{\repo}{\texttt{ProveIt}}
\newcommand{\pin}{\texttt{b2341af5370cd8a43f72ebb2341c0d7d7395b634}}
\newtheorem{theorem}{Theorem}[section]
\newtheorem{lemma}[theorem]{Lemma}
\newtheorem{proposition}[theorem]{Proposition}
\newtheorem{corollary}[theorem]{Corollary}
\theoremstyle{definition}
\newtheorem{definition}[theorem]{Definition}
\newtheorem{example}[theorem]{Example}
\theoremstyle{remark}
\newtheorem{remark}[theorem]{Remark}
\setlist{itemsep=3pt,topsep=5pt}
\setlength{\emergencystretch}{3em}
\setlength{\parskip}{3pt}
\pagestyle{fancy}
\fancyhf{}
\fancyhead[L]{\small Continuous mathematics in iPRA}
\fancyhead[R]{\small 20 November 2026}
\fancyfoot[C]{\thepage}
\lstset{basicstyle=\ttfamily\footnotesize,breaklines=true,columns=fullflexible,keepspaces=true,frame=single,rulecolor=\color{black!20},backgroundcolor=\color{black!2},showstringspaces=false}
\begin{document}
\begin{titlepage}
\vspace*{14mm}
{\Large\color{MidnightBlue} Constructive mathematics / Research extension\par}
\vspace{14mm}
{\huge\bfseries Continuous mathematics in iPRA\par}
\vspace{6mm}
{\LARGE Measures, stochastic operators,\par and homological transfer\par}
\vspace{12mm}
{\large Prepared for Vladimir Reshetnikov\par}
{\large 20 November 2026\par}
\vspace{12mm}
\begin{minipage}{0.94\textwidth}
\textbf{Purpose.} Extend the repository's finite-presentation framework with a
measure-sensitive calculus that preserves the distinction between uniform,
integral, and operator error. The applications include a corrected constructive
Vitali theorem with an explicit limit constructor, conditional expectations,
probability-preserving operators, integrated continuous formulas, holomorphic
parameter integrals, and quantitative homological perturbation.

\medskip
\textbf{Proof status.} The report gives human-readable constructive proofs and
explicit arithmetic compilation obligations. The accompanying standalone Lean
project checks finite quantitative kernels. Its new theorems are host Lean
theorems; the full analytic results have not been exported as new proof values
of type \texttt{Prf PRAAx}. The repository's existing object-theory arithmetic
is identified precisely for reuse. No unrestricted interpretation of analysis
in first-order arithmetic is claimed.
\end{minipage}
\vfill
\textbf{Reviewed repository snapshot}\par
\smallskip
\url{https://github.com/VladimirReshetnikov/ProveIt}\par
\smallskip
{\small\pin\par}
\end{titlepage}
\begin{abstract}
Finite observations become substantially more useful when their error is
measured in the topology consumed by the next construction. We develop this
principle for probability-algebra presentations and normed complexes. Stable
clipping replaces discontinuous event tests and repairs a false inequality in
an existing Vitali argument. Quantitative uniform integrability then upgrades
a Cauchy name in measure to a constructed $L^1$ name, without assuming an
integrable limit in advance. Finite conditional averages extend by contraction
and inherit explicit approximation rates. A category of presented
probability-preserving operators composes through strong action data, avoiding
an unjustified operator-norm approximation of the identity. Finite nets and
probability quadratures support a mixed infimum/supremum/integration logic with
syntax-directed error bounds. Integration of $L^1$-valued power series gives
effective holomorphic parameter integrals. Finally, a signed homological
perturbation construction uses either a supplied nilpotence bound or a
certified norm gap, with explicit residuals and Neumann tails. All numerical
schedules are primitive recursive in the stated data. We separate these
mathematical arguments, checked Lean kernels, and the remaining first-order
proof compilation work.
\end{abstract}
\tableofcontents
\clearpage
\section{The extension and its exact scope}\label{sec:scope}

The current repository is already far beyond the proposal that continuous
objects should be replaced by finite certificates. Its November 20 packets
include quantitative completion, finite nets of actual Lipschitz functions,
continuous model semantics, holomorphic gluing, coercive operators, and
topological descent. Its \texttt{Derivations} tree contains substantial
\emph{object-theory} arithmetic, including rational sums, weighted error
budgets, and finite conditional-expectation identities
\cite{repo-overview,repo-derivations,repo-abstract,repo-continuous-two}.
The overview's older packet counts should therefore not be treated as a
complete inventory. We inspected the pinned source tree directly.

This report extends that work in a direction where the previous abstractions
need additional structure: a small error in mean, a small uniform error, and
a small operator error support different mathematical operations. Treating
them interchangeably can turn a correct finite computation into an invalid
continuous conclusion. Conversely, demanding uniform error everywhere can
make useful constructions needlessly expensive or impossible.

\subsection{What is reused and what is added}

\begin{longtable}{@{}p{0.23\textwidth}p{0.34\textwidth}p{0.36\textwidth}@{}}
\toprule
Area & Existing material inspected & Contribution here\\
\midrule
\endhead
Arithmetic and proof composition & \texttt{PRA.qfInd}, rational arithmetic and finite sums, finite Fubini, proof-carrying error DAGs & A concrete route for clipping, mean-error, and kernel estimates through those interfaces; the arithmetic infrastructure is credited rather than reintroduced as new.\\
Completion and topology & Regular names, dense presentations, compact nets, coherent witnesses & Norm-sensitive extension: $L^1$, convergence in measure, strong operator action, and exact chain identities have distinct interfaces.\\
Expectation & Finite natural weights, common denominators, conditional cell regrouping, a dominated truncation bound & Stable clipping-Vitali estimate, an $L^1$ limit constructor from Cauchy-in-measure/UI data, and completion of finite conditional averages.\\
Category theory & Presented setoids, completion functors, finite diagrams and gluing & A composable category of probability-preserving operators; its tensor and identity are justified in the appropriate norms.\\
Model theory & Quantitative transfer for compact continuous formulas with infimum and supremum & Integration binders and mixed error budgets; preservation of binder order and explicit limits of $L^1$ transfer.\\
Complex analysis & Certified power-series and derivative tails & Integration of $L^1$-valued series, differentiation under the integral, and a complete rational example.\\
Algebraic topology & Finite nerve and chain comparisons; a separate $\mathbb F_2$ perturbation calculation & Signed normed perturbation with a finite nilpotence clock or a Neumann gap, quantitative residuals, and stability.\\
\bottomrule
\end{longtable}

The theory of Markov categories, probability algebras, quantitative convergence,
and the homological perturbation lemma are established subjects
\cite{fritz,berenstein-henson,neri-pischke,avigad-dean-rute,crainic2004}.
The contribution is their explicit, compatible use within this repository's
first-order presentation discipline, together with the correction and
additional quantitative constructions below. We do not claim a new classical
Vitali theorem or a new homological perturbation lemma.

\subsection{A specific mathematical repair}

The Vitali proof in
\texttt{expected\_value\_ipra\_2026-11-16}, Section 12, uses the pointwise claim
\begin{equation}\label{eq:false-old-vitali}
 |X-Y|\le |X-Y|\mathbf1_{|X|\le K,\,|Y|\le K}
       +|X|\mathbf1_{|X|>K}+|Y|\mathbf1_{|Y|>K}.
\end{equation}
On the one-point probability space, set $K=1$, $X=2$, $Y=-1$.
The two sides are $3$ and $2$, respectively. This is a defect in the
displayed proof, not a counterexample to the stated Vitali conclusion
\cite{repo-expectation-packet}. The separate checked
\texttt{Probability.TruncationEstimate} uses a common dominating function
and a different, valid three-region estimate; it does not assert
\eqref{eq:false-old-vitali} \cite{repo-truncation}.

The replacement uses
\[
 c_K(t)=\max(-K,\min(t,K)),\quad
 t_K(X)=\int (|X|-K)_+,\quad
 \rho(X,Y)=\int\min(1,|X-Y|).
\]
For $K\ge1$ we prove
\begin{equation}\label{eq:main-vitali-intro}
 \norm{X-Y}_1\le t_K(X)+t_K(Y)+2K\rho(X,Y).
\end{equation}
All three quantities on the right are stable under $L^1$ approximation.
Unlike a hard threshold probability, they require no decision about equality
at a boundary. The theorem subsequently constructs the limit from pairwise
estimates and explicit moduli; it does not presuppose that the desired limit
has already been represented in $L^1$.

\subsection{How to read the claims}

There are three distinct deliverables. First, this report proves statements
about supplied presentations and quantitative hypotheses. Second, the Lean
project proves selected finite arithmetic and logical kernels. Third, the
object-language plan specifies the numerical functions, quantifier-free
invariants, existing interfaces, and missing proof values needed to export
the new results to \texttt{Prf PRAAx}. A checked host theorem does not itself
complete that third task.

The infinite notation is a readable account of numerical constructions.
The iPRA content is always a theorem about finite observations at arbitrary
natural-number precisions. Each evaluator, refinement operation, tail
modulus, and clock used in a uniform claim must be supplied as a fixed PR
definition, or be implemented by a bounded interpreter with a separately
proved clock. Neither arbitrary real functions nor arbitrary measures are
being quantified over as objects of first-order iPRA.

The ZIP includes editable sources, a compiled PDF, the complete standalone
Lean sources, exact-rational reference programs, an immutable source map,
and a claim-status ledger. No remote repository files were changed.

\section{A calculus that records the norm of each observation}\label{sec:foundation}

\subsection{Arithmetic, coding, and the meaning of a presentation}

We use intuitionistic first-order logic with equality, the natural numbers,
a symbol for every fixed primitive-recursive definition, its defining
equations, successor injectivity and nonzeroness, and induction for
quantifier-free formulas. This is the repository's ordinary
\texttt{PRAAx} presentation \cite{repo-pr}. No induction over arbitrary
quantified statements is added.

A signed integer is a difference of two naturals. A rational is an integer
numerator and a positive natural denominator; it can be normalized by the
PR gcd algorithm but normalization is not needed for order. Equality of
$a/b$ and $c/d$ means $ad=bc$, and order means $ad\le bc$, with $b,d>0$.
A finite list is coded by the repository's pairing and sequence operations.
Every access is explicitly guarded by its stored length. Finite vectors,
matrices, Boolean tables, and polynomials are lists with dimension and
well-formedness checks. None requires a new arithmetic sort.

For clarity, a map displayed as $a:\N\to\Q$ below denotes a \emph{fixed
program schema} and its numerical parameters. It does not denote an
arbitrary function variable in the language. A uniform statement over
program codes uses a bounded evaluator $\operatorname{run}(e,x,t)$, defined
by primitive recursion on $t$, together with a supplied PR clock and its
termination proof. There is no total PR universal evaluator for all PR
programs: if $U(e,x)$ were one, $U(n,n)+1$ would be a PR function and its own
index would give a contradiction. This elementary diagonal argument is a
useful guard against an otherwise tempting but invalid encoding shortcut.

\begin{definition}[A quantitative presentation]\label{def:presentation}
A presentation supplies a collection of finite codes, a PR admissibility
test, a PR common-refinement operation where needed, and an exactly
computable rational pseudodistance $d$ on finite codes after refinement.
The pseudometric laws and refinement invariance are proved from the
finite definitions. A regular name is a fixed PR sampler $a(n)$ satisfying
\begin{equation}\label{eq:regular-name}
 d(a(n),a(m))\le2^{-n}+2^{-m}\qquad(n,m\in\N).
\end{equation}
Two names are equivalent when their distance tends to zero with the given
regular estimates; equivalently, after a harmless precision shift, their
cross-distances are bounded by $2^{-n}+2^{-m}$ for all $n,m$.
\end{definition}

The usual metric completion is one semantic interpretation of these names.
The arithmetic implementation keeps the samplers and the estimates and
need not construct a quotient set. A fixed sampler's regularity is a
first-order formula universally quantifying $n,m$, with a quantifier-free
matrix. Subsequent arguments may use that formula as an assumption; their
induction invariants remain quantifier-free.

If finite distances themselves have certified rational upper and lower
approximations, the same construction works with extra guard bits. The
applications here deliberately use exact rational norms on simple real
functions and finite matrices whenever possible. Complex quantities can be
represented by two real codes and bounded by the rational norm
$|a+ib|_\square=|a|+|b|$; Euclidean-radius guards can instead be checked by
squaring nonnegative rational bounds. These equivalent norm presentations
are connected by explicit inequalities, not identified literally.

\begin{lemma}[Diagonal completion with a supplied clock]\label{lem:diagonal}
Suppose $x_j$ are presented points with jointly PR names, and $N$ is a
nondecreasing PR function such that
\[
 j,\ell\ge N(p)\quad\Longrightarrow\quad d(x_j,x_\ell)\le2^{-p}.
\]
Let $b(p)$ be a finite approximation to $x_{N(p+2)}$ within $2^{-p-2}$.
Then $b$ is a regular name. Its limit $x$ satisfies
$d(x_j,x)\le2^{-p}$ for $j\ge N(p)$.
\end{lemma}
\begin{proof}
For $p\le q$, the common Cauchy estimate gives
\[
 d(b(p),b(q))\le2^{-p-2}+2^{-p-2}+2^{-q-2}
 \le2^{-p}+2^{-q}.
\]
For fixed $j\ge N(p)$, compare $x_j$ with $b(q)$ for $q$ large enough that
$N(q+2)\ge N(p)$. The bound is $2^{-p}+2^{-q-2}$; the named limit removes
the final term. This last sentence is arithmetic shorthand for all its
rational-precision instances. To prove any strict enlargement of $2^{-p}$,
choose a larger explicit $q$; no choice of an unbounded limit index occurs.
Replacing $p$ by $p+1$ supplies a direct weak $2^{-p}$ output bound if a
one-sided name convention is desired. All samplers are compositions of
the supplied PR programs.
\end{proof}

\subsection{Norms are part of the interface}

For a finite probability vector $w=(w_i)_{i<n}$, $w_i\ge0$,
$\sum_iw_i=1$, define
\[
 I_w(f)=\sum_{i<n}w_i f_i,\qquad
 \norm f_{1,w}=\sum_iw_i|f_i|,\qquad
 \norm f_\infty=\max_i|f_i|.
\]
The maximum over an empty list is defined to be zero, although a normalized
probability vector cannot have empty support. Refinement repeats values
on child cells and splits mass, preserving the first two expressions.
For a real matrix use the row-sum norm
$\norm A_\infty=\max_i\sum_j|A_{ij}|$, including dimension-zero cases with
the same convention. All these are PR rational functions.

An operation's contract names both the input and output error. Integration
is $L^1\to|\cdot|$ nonexpansive; taking a supremum is
$\|\cdot\|_\infty\to|\cdot|$ nonexpansive; multiplication is
$L^\infty\times L^1\to L^1$ continuous with a supplied bounded-factor
constant. These distinctions are logically significant.

\begin{proposition}[Safe elementary transfers]\label{prop:safe-transfers}
On a finite probability algebra, and hence on the indicated completions,
\begin{align}
 |I(f)-I(g)|&\le\norm{f-g}_1,\label{eq:int-lipschitz}\\
 \norm{fg-\widetilde f\widetilde g}_1
 &\le B\norm{g-\widetilde g}_1+C\norm{f-\widetilde f}_1,
 \label{eq:product-lipschitz}
\end{align}
provided $|f|\le B$ and $|\widetilde g|\le C$ almost everywhere.
Also $|\sup f-\sup g|\le\norm{f-g}_\infty$ when those suprema are named.
\end{proposition}
\begin{proof}
The first inequality is the triangle inequality for a finite weighted sum.
For the second write $fg-\widetilde f\widetilde g
=f(g-\widetilde g)+(f-\widetilde f)\widetilde g$ and apply the pointwise
triangle inequality before summation. For a supremum, $f(x)\le
g(x)+\norm{f-g}_\infty$ gives one direction, and interchanging $f,g$ gives
the other. Completion preserves inequalities between uniformly continuous
operations by comparing finite approximants with arbitrarily small explicit
errors. The bounded-factor hypotheses in the product estimate must be
preserved or re-certified; $L^1$ names alone do not supply them.
\end{proof}

There is no general transfer from $L^1$ error to uniform error. On $[0,1]$,
$f_n=\mathbf1_{[0,2^{-n}]}$ has $\norm{f_n}_1=2^{-n}$ and
$\norm{f_n}_\infty=1$. Even continuous tent functions give the same
obstruction. Likewise, finite-support atomic measures do not approximate
Lebesgue measure in total variation: their finite support has mass one
for the atomic measure and zero for Lebesgue measure. Finite cell
\emph{observations} of a measure and atomic approximations for Lipschitz
test functions are therefore separate representations. The later category
and model-theory sections use the representation appropriate to each claim.

\subsection{Weighted error and efficient demand propagation}

\begin{lemma}[Mass and value perturbation]\label{lem:weighted}
Let $w_i\ge0$, let $v_i$ be rational, and suppose
$|a_i-b_i|\le e_i$ and $|b_i|\le B$. Then
\begin{equation}\label{eq:weighted}
 \left|\sum_i w_i a_i-\sum_i v_i b_i\right|
 \le\sum_iw_i e_i+B\sum_i|w_i-v_i|.
\end{equation}
If $w=v$ is a probability vector and all $e_i\le e$, the bound is $e$.
\end{lemma}
\begin{proof}
Subtract and add $\sum_iw_ib_i$. The difference is
$\sum_iw_i(a_i-b_i)+\sum_i(w_i-v_i)b_i$.
Apply the triangle inequality to each summand, then use the two supplied
pointwise bounds. This proof works for zero weights and does not divide by
any of them. It also makes explicit which vector weights each local error.
\end{proof}

The repository already has an internal rational error-DAG producer and a
uniform precision schedule \cite{repo-dag,repo-dag-schedule}. A useful
extension is to retain the actual probability coefficients and the norm
type of every edge. The following standard finite calculation explains why.

\begin{proposition}[Backward influence coefficients]\label{prop:adjoint}
For a topologically ordered finite computation suppose its errors satisfy
\[
 E_v\le\delta_v+\sum_{u<v}a_{vu}E_u,\qquad a_{vu}\ge0.
\]
For a selected output $o$, initialize $\lambda_o=1$ and other coefficients
to zero. Process vertices in decreasing order, adding
$\lambda_v a_{vu}$ to $\lambda_u$ along every incoming edge. Then
\[
 E_o\le\sum_v\lambda_v\delta_v.
\]
Consequently any nonnegative local budgets with
$\sum_v\lambda_v\delta_v\le2^{-p}$ certify the output precision.
\end{proposition}
\begin{proof}
Begin with $E_o$. Substitution of the inequality for a vertex replaces its
coefficient times $E_v$ by that coefficient times $\delta_v$ plus the
weighted predecessor errors. Nonnegativity preserves the direction of the
inequality. After each vertex has been processed once, no $E_v$ remains.
The accumulated coefficient of each $\delta_v$ is exactly the value
computed by the reverse pass. This proves the displayed bound, including
graphs with shared subexpressions. A bounded reverse loop and rational
addition/multiplication implement the pass by primitive recursion.
\end{proof}

For integration, $a_{vu}$ are the actual masses. A term of mass $2^{-20}$
can tolerate much larger absolute error than one of mass $1/2$.
Zero-influence terms need no numerical refinement once their range and
well-formedness are established. This is a mathematical error reduction,
not a claim that every resulting algorithm is polynomial-time. Tail
moduli, matrix dimensions, and compact-net sizes may still be enormous.
Max-type rules can be kept as maxima or safely overestimated by sums;
using a sum changes the cost bound but not correctness.

\subsection{Why the finite arguments fit quantifier-free induction}

For \eqref{eq:weighted}, fix a common positive denominator for all finite
input values. After clearing it, let $S_j,T_j,U_j,V_j$ be the relevant
integer prefix sums. A convenient induction invariant is
\[
 -U_j-BV_j\le S_j-T_j\le U_j+BV_j,
\]
together with the stored positivity conditions. This is quantifier-free
in the codes and $j$. At the successor step add the single new weighted
inequality. The zero case is literal equality. The normalization step uses
positive-denominator order preservation already available in the
repository's rational arithmetic. The variable-length sum is thus an
internal numerical function, not an externally expanded expression for
one fixed dimension.

A bounded universal condition is represented by a PR conjunction fold:
$\operatorname{all}_P(0)=1$ and
$\operatorname{all}_P(j+1)=\operatorname{and}
(\operatorname{all}_P(j),P(j))$, where $P$ is a decidable numerical test.
Proofs about a whole input table can use the equality
$\operatorname{all}_P(n)=1$ as a quantifier-free premise and derive its
individual bounded instances. In every main construction below we give
either the corresponding prefix invariant or the finite algebraic identity
that replaces it. There is no rule saying that an arbitrary true numerical
claim is automatically provable in iPRA.

\section{Measure and integration from coherent finite observations}
\label{sec:measure}

The repository already has an effective expectation theory, a dominated
convergence statement, and a Vitali statement
\cite{repo-expectation-packet}.  Its current object-language development is
substantial: \texttt{FiniteKernel.pr\_evaluator} makes cross-multiplied
expectation comparisons quantifier-free; \texttt{TowerCommonDen.regroup}
proves the numerator regrouping used by the tower property; and
\texttt{TruncationEstimate.truncation} proves a finite inequality with a
common dominator
\cite{repo-finite-kernel,repo-tower-common-den,repo-truncation}.
The new work here concerns the passage from these finite identities to
coherent measure-theoretic objects.  It supplies a stable convergence
observable, constructs an integrable limit rather than presupposing one,
extracts amplitude cutoffs from an existing integrable name, and extends
finite conditional averages with an explicit approximation rate.

There is also a specific local correction.  The earlier packet's human
Vitali proof uses
\[
 |x-y|\le |x-y|\mathbf1_{|x|\le K,\ |y|\le K}
       +|x|\mathbf1_{|x|>K}+|y|\mathbf1_{|y|>K}.
\]
For $K>0$, the values $x=2K$, $y=-K$ give $3K\le2K$, so that displayed
inequality is false.  This does not affect the different, common-dominator
inequality in the checked \texttt{TruncationEstimate} module.  The clipping
estimate proved below repairs the Vitali argument and also avoids computing
hard threshold events from approximations to real values.

Throughout this section put $\delta_k=2^{-k}$.  Every input evaluator,
refinement map and modulus mentioned as effective is a \emph{fixed}
primitive-recursive program. For an iPRA export, its displayed numerical
contracts must be supplied as object-theory derivations or retained as
explicit premises. Numerical problem parameters may vary. This is a
schema about fixed programs, not a claim that arbitrary codes for total
programs have a primitive-recursive universal evaluator.

\subsection{A common probability algebra and its integrable completion}

\begin{definition}[Coherent rational probability presentation]
For each $n\in\N$, let $A_n=\{0,\ldots,a(n)-1\}$ be a nonempty finite set
of atoms, with rational weights $w_{n,i}\ge0$ summing to one.  For $m\le n$
there is a map $\pi_{m,n}:A_n\to A_m$, satisfying
\begin{align}
 \pi_{n,n}(i)&=i,&
 \pi_{\ell,n}&=\pi_{\ell,m}\circ\pi_{m,n}\quad(\ell\le m\le n),
 \label{eq:m-refinement}\\
 w_{m,i}&=\sum_{\substack{j<a(n)\\\pi_{m,n}(j)=i}}w_{n,j}.
 \label{eq:m-mass}
\end{align}
The cardinalities, weights and maps have fixed PR evaluators.  Their range,
normalization and compatibility equations are part of the presentation.
\end{definition}

A simple function is a code $s=(n,(s_i)_{i<a(n)})$, with rational values.
To compare $s$ at level $n$ with $t$ at level $m$, lift both lists to
$r=\max(m,n)$ using $\pi_{n,r}$ and $\pi_{m,r}$.  On that common level set
\begin{equation}
 I(s)=\sum_iw_i s_i,\qquad
 \|s\|_1=\sum_iw_i|s_i|,\qquad
 d_1(s,t)=\sum_iw_i|s_i-t_i|.
 \label{eq:m-simple}
\end{equation}
Equation~\eqref{eq:m-mass} proves that lifting again changes none of these
values.  It also proves independence of the chosen common refinement.
Finite addition, scalar multiplication, minimum, maximum and absolute value
are computed on a common level.  Zero-weight cells are retained in codes
but have zero distance.  Thus this is a rational vector lattice with a
pseudometric, rather than a demand for decidable equality of functions.

The coherent tower matters.  Marginal probability laws alone do not
determine a common $L^1$ geometry: two fair signs can be equal or opposite,
giving respectively zero or two as the expectation of their absolute
difference.  The maps \eqref{eq:m-refinement} choose the common probability
algebra explicitly and keep every operation on that algebra.

This presentation has concrete continuous models.  For dyadic intervals
in $[0,1)$, take $a(n)=2^n$, $w_{n,i}=2^{-n}$ and
$\pi_{m,n}(i)=\lfloor i/2^{n-m}\rfloor$.  More generally, every tower above
has an explicit interval realization.  Give each atom of $A_0$ a half-open
interval of its stated length.  Inside the interval for $i\in A_n$, put
the children $j$ with $\pi_{n,n+1}(j)=i$ in increasing order, with lengths
$w_{n+1,j}$.  The left endpoint of child $j$ is the parent's left endpoint
plus the sum of the weights of its earlier siblings.  Compatibility makes
these children fill the parent interval exactly.  Zero weights give empty
intervals.  This recursive rational construction realizes all finite
integrals as ordinary step-function integrals.  No infinite branch is
selected.  In a usual semantic $L^1$ space the completion below is precisely
the closure of these step functions; asserting density in a larger space
is a separate adequacy statement.  The interval realization alone neither
codes all measurable subsets nor supplies a rate for arbitrary countable
unions.

\begin{definition}[Regular integrable name]
A regular $L^1$ name is a fixed PR sequence of simple-function codes
$s_k$ satisfying
\begin{equation}
 d_1(s_k,s_\ell)\le\delta_k+\delta_\ell
 \qquad(k,\ell\in\N).
 \label{eq:m-regular}
\end{equation}
Its semantic value $f$ is the corresponding completion element; in
particular $\|f-s_k\|_1\le\delta_k$.  Equivalent names denote the same
completion element.  Internally one works with
\eqref{eq:m-regular} and finite comparisons of the actual codes, not with
a new sort of arbitrary measurable functions.
\end{definition}

For example, the distance between two represented elements is itself a
regular scalar name: evaluate $d_1(s_{k+1},t_{k+1})$; the shift absorbs
both input errors.  Addition similarly uses $s_{k+1}+t_{k+1}$ at output
precision $k$.  Equality means that every requested
distance observation can be made arbitrarily small.  It is an arithmetical
formula, not an algorithm which tests exact real equality.

\begin{lemma}[Stable extension of finite observables]
\label{lem:m-extension}
The integral and norm extend to integrable names, with
\begin{equation}
 |I(f)-I(g)|\le\|f-g\|_1,\qquad
 \bigl|\|f\|_1-\|g\|_1\bigr|\le\|f-g\|_1.
 \label{eq:m-integral-stable}
\end{equation}
Any finite operation $T$ satisfying
$d_1(Ts,Tt)\le d_1(s,t)$ extends by the same precision schedule.  This
applies to $s\mapsto |s|$, $s\mapsto s^+=\max(s,0)$, and clipping to a
fixed rational interval.  Finite identities between such operations
remain true on the completion.
\end{lemma}

\begin{proof}
The finite triangle inequality gives
$|\sum_iw_i(s_i-t_i)|\le\sum_iw_i|s_i-t_i|$.  The reverse triangle
inequality gives the second assertion in \eqref{eq:m-integral-stable}.
Consequently the finite values on a regular name form regular scalar
names.  A nonexpansive $T$ carries \eqref{eq:m-regular} to the identical
inequality for $Ts_k$.  If two input names represent the same element,
their image names do also, by this inequality.  A finite identity is an
identity at every common approximation level; each side's error then tends
to zero with its displayed schedule.  This proves existence,
well-definedness and preservation, without choosing representatives at
individual sample points.
\end{proof}

The completed order is $f\le g$ when $\|(f-g)^+\|_1=0$.  It is stable
under limits by the same Lipschitz argument.  Multiplication requires its
own contract: a bounded multiplier $|h|\le M$ satisfies
$\|hf-hg\|_1\le M\|f-g\|_1$.  Products of two unrestricted $L^1$
elements are not part of this interface.

Finite events are subsets of a level's atom set.  Their metric is
$d_\mu(B,C)=\mu(B\mathbin\triangle C)$; their indicator functions embed
isometrically in the simple $L^1$ space.  Complement preserves this metric,
and union and intersection change by at most the sum of the input errors.
Thus regular event names form an effective measure algebra.  For a
uniformly PR-presented sequence of event names, a countable union can be
constructed when its finite partial unions $B^{<n}$ have a PR continuity
modulus
\[
 m\ge n\ge R(k)\quad\Longrightarrow\quad
 \mu(B^{<m}\setminus B^{<n})\le\delta_k.
\]
Replace $R$ by its bounded monotone envelope if needed.  To make the
diagonal explicit, take $N=R(k+2)$, approximate each of the
first $N$ events to error
$2^{-k-2-\lceil\log_2(N+1)\rceil}$, and take their finite union.  Its
total sampling error is at most $\delta_{k+2}$.  Together with the
displayed continuity bound this gives a regular output name by the same
three-error estimate as \eqref{eq:m-output-name} below.  All logarithms
here mean bounded integer binary-length calculations.  For disjoint
events this premise is exactly a supplied tail estimate for their mass
series; finite additivity and continuity of $\mu$ give countable
additivity for this represented union.  An arbitrary sequence of event
codes does not itself supply the displayed modulus.

\subsection{Stable observations of convergence in measure}

For $K\ge0$, define the scalar clipping operation and excess by
\[
 c_K(x)=\max(-K,\min(K,x)),\qquad
 p_K(x)=(|x|-K)^+.
\]
For integrable elements set
\begin{equation}
 T_K(f)=I(p_K(f))=\|f-c_K(f)\|_1,
 \qquad
 \rho(f,g)=I\bigl(\min(1,|f-g|)\bigr).
 \label{eq:m-stable-observers}
\end{equation}
Both expressions have evaluators inherited from finite rational lists.
The function $p_K$ is 1-Lipschitz, as is $x\mapsto\min(1,|x|)$.
Consequently
\begin{equation}
 |T_K(f)-T_K(g)|\le\|f-g\|_1,\quad
 |\rho(f,g)-\rho(f',g')|
 \le\|f-f'\|_1+\|g-g'\|_1.
 \label{eq:m-observer-lipschitz}
\end{equation}
These bounds are the reason finite approximation errors can be transported
through the observers.

\begin{lemma}[Clipping estimate]
\label{lem:m-clipping}
For $K\ge1$ and integrable $f,g$,
\begin{equation}
 \|f-g\|_1\le T_K(f)+T_K(g)+2K\rho(f,g).
 \label{eq:m-clipping}
\end{equation}
Also $\rho(f,g)\le\|f-g\|_1$, and $\rho$ satisfies the triangle inequality.
\end{lemma}

\begin{proof}
For rational $x,y$, the clipping map is nondecreasing and 1-Lipschitz,
has range $[-K,K]$, and $|x-c_K(x)|=p_K(x)$.  Hence
\begin{align*}
 |x-y|&\le p_K(x)+p_K(y)+|c_K(x)-c_K(y)|,\\
 |c_K(x)-c_K(y)|&\le\min(|x-y|,2K)
                         \le2K\min(1,|x-y|).
\end{align*}
For the last inequality, split the rational cases $|x-y|\le1$ and
$|x-y|\ge1$.  Multiply by the nonnegative weights and add.  This proves
\eqref{eq:m-clipping} on every finite common refinement.  For names,
replace each argument by a simple approximation and use
\eqref{eq:m-observer-lipschitz}; the total error is at most a fixed rational
multiple of the two approximation errors, so it has an explicit PR
precision schedule.  This establishes the same inequality in the
completion.  Finally,
$\min(1,a+b)\le\min(1,a)+\min(1,b)$ for $a,b\ge0$, which gives the
triangle inequality for $\rho$ after the scalar triangle inequality.
\end{proof}

The metric $\rho$ expresses convergence in measure without a threshold
decision.  In an ordinary probability-space interpretation, for
$0<\eta\le1$ the finite inequalities give
\begin{equation}
 \eta\,\mu(|f-g|>\eta)\le\rho(f,g)
 \le\eta+\mu(|f-g|>\eta).
 \label{eq:m-probability-comparison}
\end{equation}
Thus a two-parameter convergence-in-probability rate converts to a
$\rho$-rate, and conversely.  In the abstract numerical presentation, the
hard event in \eqref{eq:m-probability-comparison} is used only when its
measure has separately been justified.  The algorithm itself evaluates
the continuous truncated distance in \eqref{eq:m-stable-observers}.
An $L^1$ approximation alone does not compute the measure of a sharp
threshold event: constants just below and above a threshold can be
arbitrarily close in $L^1$ while their threshold-event probabilities are
zero and one.

One can also complete the simple space using $\rho$.  Call that completion
$L^0$ for this presentation.  Since $\rho\le d_1$, there is a canonical
map $L^1\to L^0$.  It is injective.  Indeed, for any fixed integrable $h$,
choose a simple $s$ with $\|h-s\|_1\le\varepsilon$ and
$K\ge\max(1,\|s\|_\infty)$.  Then $T_K(h)\le\varepsilon$ by
\eqref{eq:m-observer-lipschitz}.  If $\rho(h,0)=0$, the clipping estimate
gives $\|h\|_1\le\varepsilon$ for every supplied precision.  The theorem
below explicitly constructs an $L^1$ preimage of an $L^0$ Cauchy limit.

\subsection{A quantitative Vitali theorem that constructs its limit}

\begin{theorem}[Uniform-integrability lifting]
\label{thm:m-vitali}
Let $f_n$ be a uniformly presented family of integrable names: a fixed PR
program $S(n,j)$ returns a simple code with
$\|S(n,j)-f_n\|_1\le\delta_j$.  Suppose that monotone PR functions
$U,R:\N\to\N$, with $U(k)\ge1$, satisfy
\begin{align}
 T_{U(k)}(f_n)&\le\delta_k &&\text{for every }n,k,
 \label{eq:m-ui}\\
 \rho(f_m,f_n)&\le\delta_k &&\text{if }m,n\ge R(k).
 \label{eq:m-rho-cauchy}
\end{align}
For an integer $b\ge1$ put
\[
 \ell(b)=\min\{j\le b: b\le2^j\}.
\]
Define the explicit precision program
\begin{equation}
 q(k)=k+2+\ell\bigl(2U(k+2)\bigr),\qquad
 N(k)=R(q(k)).
 \label{eq:m-vitali-rate}
\end{equation}
Then there is a regular integrable name $f$, produced by a PR program,
such that
\begin{equation}
 n\ge N(k)\quad\Longrightarrow\quad
 \|f_n-f\|_1\le\tfrac34\delta_k.
 \label{eq:m-vitali-convergence}
\end{equation}
In particular the integrals converge with this rate.  The limit satisfies
$T_{U(k)}(f)\le\delta_k$; no integrable limit or tail bound for a limit is
an input assumption.
\end{theorem}

\begin{proof}
The bounded definition of $\ell$ is total because $2^b\ge b$; this finite
inequality is proved by quantifier-free induction.  Consequently
$2U(k+2)2^{-\ell(2U(k+2))}\le1$.  If $m,n\ge N(k)$, apply
Lemma~\ref{lem:m-clipping} with $K=U(k+2)$ and use the two input contracts:
\begin{align}
 \|f_m-f_n\|_1
 &\le2\delta_{k+2}+2U(k+2)\delta_{q(k)}
 \le3\delta_{k+2}=\tfrac34\delta_k.
 \label{eq:m-pair-cauchy}
\end{align}
This derives an $L^1$ Cauchy modulus before introducing any limit.
Monotonicity of $U,R,\ell$ makes $N$ monotone.  Nonmonotone supplied
moduli can first be replaced by their bounded maxima over initial
segments; both contracts remain valid.

Now define an actual simple-code output program by
\begin{equation}
 b_k=S(N(k+2),k+2).
 \label{eq:m-output-name}
\end{equation}
If $j\ge k$, both family indices in $b_k,b_j$ are at least $N(k+2)$.
Their represented values have distance at most $\delta_{k+2}$ by
\eqref{eq:m-pair-cauchy}.  Adding the two simple-approximation errors gives
\[
 d_1(b_k,b_j)\le2\delta_{k+2}+\delta_{j+2}
              \le\tfrac34\delta_k\le\delta_k+\delta_j.
\]
Thus $b_k$ is a regular name, and defines $f$.  For $n\ge N(k)$ and
all sufficiently large $j$, \eqref{eq:m-pair-cauchy} gives
\[
 \|f_n-b_j\|_1\le\tfrac34\delta_k+\delta_{j+2}.
\]
The definition of completion proves \eqref{eq:m-vitali-convergence}.
Equivalently, each rational observation is obtained by choosing $j$ large
enough for its specified remaining error; there is no search for a
convergence stage.

For each fixed $k$, continuity of the tail observable yields
\[
 T_{U(k)}(f)\le T_{U(k)}(f_n)+\|f-f_n\|_1
             \le\delta_k+\|f-f_n\|_1.
\]
Substitute $n=N(j)$ and increase the requested numerical precision $j$.
This proves $T_{U(k)}(f)\le\delta_k$.  Only the explicit Lipschitz bound
for $T_K$ is used; neither Fatou's lemma nor a separately assumed
uniform-integrability property of an unknown limit is needed.
Finally \eqref{eq:m-integral-stable} gives the same estimate for
$|I(f_n)-I(f)|$.
\end{proof}

This is a quantitative inverse to the inclusion $L^1\to L^0$ on a family
equipped with $U$.  In particular, if an independently presented
$L^0$ element is the $\rho$-limit of $f_n$, it equals the image of the
constructed integrable $f$, by the triangle inequality.  The conclusion
does not assume that the independently presented limit was already
integrable.

\begin{proposition}[A named dominator supplies its own cutoff]
\label{prop:m-domination}
Suppose $g\ge0$ has a regular integrable name $g_j$, and $|f_n|\le g$ for
all $n$ in the completed lattice order.  Define
\begin{equation}
 U(k)=\max\left(1,
       \max_{j\le k}\left\lceil\|g_j\|_\infty\right\rceil\right).
 \label{eq:m-derived-ui}
\end{equation}
Then \eqref{eq:m-ui} holds.  Thus, together with the $\rho$ Cauchy rate,
these hypotheses suffice for Theorem~\ref{thm:m-vitali}; a separate
amplitude-tail oracle for $g$ is unnecessary.
\end{proposition}

\begin{proof}
For $K\ge\|g_k\|_\infty$, $p_K(g_k)=0$.  By
\eqref{eq:m-observer-lipschitz},
$T_K(g)\le\|g-g_k\|_1\le\delta_k$.
The scalar inequality
\[
 (|x|-K)^+\le (y-K)^+ + (|x|-y)^+
\]
shows, after integration and stable passage to names, that $|f_n|\le g$
implies $T_K(f_n)\le T_K(g)$.  Apply this with $K=U(k)$.
The maximum and ceiling in \eqref{eq:m-derived-ui} inspect finite rational
tables and are PR.  They do not compute an essential supremum of $g$.
The constructed limit remains dominated: the norm of
$(|f|-g)^+$ is at most $\|f-f_n\|_1$ and hence is zero.
\end{proof}

The excess convention for $T_K$ is equivalent, with explicit constants,
to the usual uniform-integrability convention when sharp tail integrals
are available.  Pointwise,
\[
 (|x|-K)^+\le |x|\mathbf1_{|x|>K},\qquad
 |x|\mathbf1_{|x|>2K}\le2(|x|-K)^+.
\]
The excess is preferable as the primitive observer because it is
Lipschitz.  General quantitative and metastable convergence work provides
broader context for making these data explicit
\cite{neri-pischke,avigad-dean-rute}; the displayed schedules and proofs here
follow directly from finite rational inequalities.

\subsection{Conditional expectation and effective filtration convergence}

For a simple function $s$ at level $m$ and a target level $n$, let
$r=\max(m,n)$ and let $\widetilde s$ be its lift to $A_r$.  Define a
rational list on $A_n$ by
\begin{equation}
 (E_ns)_i=
 \begin{cases}
 \displaystyle\frac{\sum_{\pi_{n,r}(j)=i}w_{r,j}\widetilde s_j}{w_{n,i}},
       &w_{n,i}>0,\\[6pt]
 0,&w_{n,i}=0.
 \end{cases}
 \label{eq:m-condexp}
\end{equation}
The branch condition is an exact rational test.  If $w_{n,i}=0$, every
descendant weight is zero, since they are nonnegative and sum to it.
Thus the chosen value zero on that cell has no effect on any $L^1$
observation.  This zero-mass argument is part of the algorithm's proof,
not an implicit division convention.

\begin{theorem}[Conditional expectation on the completion]
\label{thm:m-filtration}
The finite operators \eqref{eq:m-condexp} extend to integrable names as
positive linear contractions.  For all $m,n$,
\begin{equation}
 E_mE_n=E_{\min(m,n)},\qquad I(E_nf)=I(f).
 \label{eq:m-tower}
\end{equation}
If $s_k$ is a regular simple name for $f$, and $s_k$ is stored at level
$r(k)$, then
\begin{equation}
 n\ge r(k+1)\quad\Longrightarrow\quad
 \|E_nf-f\|_1\le\delta_k.
 \label{eq:m-upward-rate}
\end{equation}
The family $(E_nf)_n$ has a uniform-integrability modulus derived directly
from $(s_k)$, without reciprocal lower bounds on nonzero atom masses.
\end{theorem}

\begin{proof}
First work on finite rational lists.  On a positive-mass cell $i$ the
weighted triangle inequality gives
\[
 w_{n,i}|(E_ns)_i-(E_nt)_i|
 \le\sum_{\pi_{n,r}(j)=i}w_{r,j}|s_j-t_j|.
\]
On a zero-mass cell both sides are zero.  Add over all $i$ and regroup
the finite double sum to obtain
\begin{equation}
 \|E_ns-E_nt\|_1\le\|s-t\|_1.
 \label{eq:m-ce-contraction}
\end{equation}
Linearity, positivity, and preservation of the integral follow from the
same cell formula; the zero-mass branch is harmless in each identity.
For $m\le n$, averaging the level-$n$ cell numerators over a level-$m$
cell counts each fine atom once, and gives exactly the level-$m$
numerator.  Division by a positive coarse weight gives $E_mE_n=E_m$;
zero coarse weight again gives zero distance.  Also $E_ms$ is already
constant on the finer level-$n$ cells, so $E_nE_ms=E_ms$ in $L^1$.
This proves the tower property.  The numerator regrouping is precisely
the reusable finite content of \texttt{TowerCommonDen.regroup}; the
normalization and zero-cell argument complete the operator statement.

By \eqref{eq:m-ce-contraction}, the sequence $E_ns_k$ is a regular name
with the original error budget, uniformly in the numerical parameter
$n$.  Lemma~\ref{lem:m-extension} therefore extends the operators and
their identities.  For $n\ge r(k)$ one has $E_ns_k=s_k$ in $L^1$.
Consequently
\[
 \|E_nf-f\|_1
 \le\|E_n(f-s_k)\|_1+\|s_k-f\|_1
 \le2\delta_k.
\]
Replace $k$ by $k+1$ to obtain \eqref{eq:m-upward-rate}.  If a monotone
rate is desired, replace $r(k+1)$ by its maximum over the first $k+2$
approximation levels.  This is a bounded PR operation.

Finally put $K_k=\max(1,\lceil\|s_k\|_\infty\rceil)$.  Each conditional
average of $s_k$ lies in $[-K_k,K_k]$, including the zero-mass branch.
For any bounded $b$ with $|b|\le K$, the scalar inequality
$(|a|-K)^+\le|a-b|$ yields
\[
 T_{K_k}(E_nf)\le\|E_nf-E_ns_k\|_1
                    \le\|f-s_k\|_1\le\delta_k.
\]
Taking bounded maxima of the $K_k$ gives the required monotone cutoff
function.  The contraction proof cancels the cell denominator against
the cell's mass; consequently no smallest positive atom weight enters
this $L^1$ precision budget.
\end{proof}

These operators are genuine conditional expectations on the represented
finite subalgebras.  For every level-$n$ event $B$, finite regrouping gives
$I(\mathbf1_BE_nf)=I(\mathbf1_Bf)$.  An element in the completed
level-$n$ range with these cell integrals is unique in $L^1$: on each
positive cell its constant value is determined by the cell integral,
and zero cells contribute no norm.  This identifies the operator with
the usual conditional expectation in the interval realization.
It does not appeal to a general Radon--Nikodym theorem for an arbitrary
sub-$\sigma$-algebra.

There is a useful distinction between observations.  To request the
conditional value at one specified positive atom to error $\eta$, the
input integral over that atom must be known to error $w_{n,i}\eta$.
That request incurs the expected reciprocal-mass precision loss.  To
request the entire conditional expectation to error $\eta$ in $L^1$,
the contraction \eqref{eq:m-ce-contraction} incurs no such loss.

\paragraph{An unbounded worked example.}
Let
$A_j=[2^{-2j-2},2^{-2j-1})\subset[0,1)$ and let $s_n$ have value $2^j$
on $A_j$ for $j<n$, and zero elsewhere.  These sets are disjoint and
$\mu(A_j)=4^{-j-1}$.  The finite geometric-sum identity gives
\[
 \|s_m-s_n\|_1
   =2^{-n-1}\bigl(1-2^{-(m-n)}\bigr)\quad(m\ge n).
\]
They therefore define an integrable, unbounded element $g$ with
$I(g)=1/2$ and $\|g-s_n\|_1=2^{-n-1}$.  Use the tower whose level $n$
exposes $A_0,\ldots,A_{n-1}$ and a single remaining atom $B_n$.  Its mass
is $(2+4^{-n})/3$.  At level $n$, conditional expectation has the original
values on the exposed annuli and the rational value
\[
 (E_ng)|_{B_n}
   =\frac{3\,2^{-n-1}}{2+4^{-n}}.
\]
Indeed, the unexposed integral is $2^{-n-1}$, as follows from the same
finite geometric identity and its explicit remainder.  On every
unexposed annulus the value of $g$ exceeds this conditional value, while
on the complement of all annuli $g=0$; that complement has mass $2/3$.
Adding these two contributions gives the exact estimate
\[
 \|E_ng-g\|_1=\frac{2^{1-n}}{2+4^{-n}}\le2^{-n}.
\]
The family $(s_n)$ is dominated by $g$.  The extracted cutoff
$U(k)=\max(1,2^{k-1})$ and the rate $R(k)=k$ meet the Vitali hypotheses,
giving the explicit sufficient rate $N(k)=2k+4$.  This example also
illustrates a concrete efficiency gain: the adapted tower has $n+1$
atoms, whereas a uniform dyadic grid resolving the first $n$ annuli has
$4^n$ cells.  The exact-arithmetic program in the archive checks the
finite tower laws, the cutoff budgets and the displayed conditional
values, including independent finite approximations to the last norm.

\begin{corollary}[Representation of a quantitatively convergent martingale]
\label{cor:m-martingale}
Let $M_n$ be uniformly PR-presented integrable elements satisfying
$E_nM_n=M_n$ and $E_mM_n=M_m$ for $m\le n$.  If they satisfy
\eqref{eq:m-ui} and \eqref{eq:m-rho-cauchy}, the name $f$ constructed in
Theorem~\ref{thm:m-vitali} satisfies $M_n=E_nf$ for every $n$.  It is the
unique such integrable element in this completion.
\end{corollary}

\begin{proof}
For fixed $n$ and every $j\ge n$, $E_nM_j=M_n$.  The contraction gives
$\|E_nf-M_n\|_1\le\|f-M_j\|_1$, which tends to zero with the explicit
Vitali schedule.  Hence $E_nf=M_n$.  If also $E_ng=M_n$ for all $n$,
Theorem~\ref{thm:m-filtration} gives $E_nf\to f$ and $E_ng\to g$,
with their respective name-dependent rates.  Comparing both limits to
one sufficiently late common $M_n$ gives $f=g$.
\end{proof}

\begin{proposition}[Why the convergence rate is additional data]
\label{prop:m-martingale-barrier}
There is a uniformly PR family of rational martingales, all bounded by
one, from which no uniform computable $L^1$ Cauchy modulus can be
extracted.
\end{proposition}

\begin{proof}
At level $n$ of the fair-bit probability tower expose bits
$\omega_1,\ldots,\omega_n$.  Delay each machine by one initial step, so
its first possible halting step is at least one.  For a machine index $e$, simulate its first
$n$ steps.  Put $M^e_n=0$ if it has not halted.  If its first halting step
is $t\le n$, put $M^e_n(\omega)=(-1)^{\omega_t}$.  This is a rational
table produced by a bounded simulation and bounded search, uniformly PR
in $e,n$ and the atom index.  If the machine has halted, the martingale
stays unchanged.  If it first halts at the next step, its new value is the
new independent fair sign, with conditional average zero.  In every
case $E_nM^e_{n+1}=M^e_n$.  The absolute bound is one, so the excess-tail
modulus is constantly one with tail zero.

Suppose a computable procedure supplied $N(e)$ such that
$\|M^e_m-M^e_n\|_1\le1/2$ for all $m,n\ge N(e)$.  Simulate $e$ through
$N(e)$ steps.  If it has not halted, it can never halt later: at a later
first halting step $t$, $M^e_{N(e)}=0$ and $|M^e_t|=1$, whose $L^1$
distance is one.  Thus this procedure would decide the halting problem.
The contradiction proves the claim.  The same argument rules out a
uniform $\rho$ Cauchy modulus, since the jump also has truncated distance
one.  It leaves intact the explicit modulus in
\eqref{eq:m-upward-rate}, whose input is already a regular integrable
name of the terminal object.
\end{proof}

\subsection{Exact internalization obligations}

The mathematical induction used above is finite: regrouping sums,
normalizing positive rational denominators, looping over atom tables,
taking bounded maxima, and proving the power-of-two budget.  Each has a
quantifier-free numerical invariant.  For example, the clipping endpoint
after multiplying by positive common denominators is the universal
closure of
\[
 K\ge1\ \Longrightarrow\
 \sum_{i<a(r)}w_i|x_i-y_i|
 \le\sum_{i<a(r)}w_ip_K(x_i)
    +\sum_{i<a(r)}w_ip_K(y_i)
    +2K\sum_{i<a(r)}w_i\min(1,|x_i-y_i|).
\]
Every symbol in this formula abbreviates a fixed PR term on numerical
codes.  The corresponding output must ultimately have type
\texttt{Prf PRAAx []} applied to the encoded formula.  A Lean proof about
$\Q$-valued functions is a useful independently checked semantic kernel,
but is not itself that derivation.

The remaining formalization path is specific.  First instantiate the
repository's rational normalization and finite sum interfaces for the
tower codes, including the zero-mass branch.  Prove the scalar clipping
and excess-Lipschitz inequalities by rational case splits; sum them using
the existing order lemmas.  Reify $U,R,\ell,q,N$ and the simple output
program \eqref{eq:m-output-name}.  For each requested finite observation,
substitute these terms into \eqref{eq:m-pair-cauchy} and the input
regularity formulas.  This produces output regularity, convergence, and
the inherited-tail observations by ordinary logical rules.  Finally
instantiate \texttt{TowerCommonDen.regroup}, prove the normalized
contraction, and obtain \eqref{eq:m-upward-rate} from the stated level
index.  None of these steps requires an induction hypothesis asserting
convergence at every precision.  Precision is a free numerical
parameter; the finite loop and all finer demands are explicit terms.

The interval model supplies semantic interpretation of the resulting
arithmetic laws.  The laws apply to the chosen completion and its
represented subalgebras.  Finite measures follow by an explicitly
positive normalization; extensions to an infinite ambient measure also
need a spatial-tail or exhaustion contract.  General pointwise
convergence, arbitrary sub-$\sigma$-algebras, and arbitrary countable
unions do not silently provide any of those effective data.

\section{Presented probability-preserving operators}\label{sec:operators}

\subsection{Finite kernels and their exact identities}

Let $(X,\mu)$ and $(Y,\nu)$ be finite probability spaces. A
probability-preserving transition is a rational matrix $P=(P_{ij})$ with
\begin{equation}\label{eq:stationarity}
 P_{ij}\ge0,\quad \sum_jP_{ij}=1,\quad
 \sum_i\mu_iP_{ij}=\nu_j.
\end{equation}
Its action on observables is contravariant:
\[
 T_P:L^1(Y,\nu)\longrightarrow L^1(X,\mu),\qquad
 (T_Pf)_i=\sum_jP_{ij}f_j.
\]
Zero-mass source rows may be chosen arbitrarily subject to row
normalization; they have no effect on the $L^1$ map. They must nevertheless
remain well-formed rows in the numerical code.

\begin{theorem}[Finite Markov-operator calculus]\label{thm:finite-markov}
The map $T_P$ is positive, linear, preserves $1$ and the integral, and
\begin{equation}\label{eq:markov-contraction}
 \norm{T_Pf-T_Pg}_{1,\mu}\le\norm{f-g}_{1,\nu}.
\end{equation}
If $Q:(Y,\nu)\to(Z,\xi)$ is another transition, its matrix composite $PQ$
satisfies \eqref{eq:stationarity}, and $T_{PQ}=T_PT_Q$.
Finite independent products are preserved by the tensor matrix
$(P\otimes Q)_{(i,k),(j,\ell)}=P_{ij}Q_{k\ell}$.
\end{theorem}
\begin{proof}
Positivity and linearity follow termwise; the row-sum law gives $T_P1=1$.
Finite Fubini and the last equation in \eqref{eq:stationarity} give
\[
 \sum_i\mu_i(T_Pf)_i
 =\sum_j\left(\sum_i\mu_iP_{ij}\right)f_j
 =\sum_j\nu_jf_j.
\]
For contraction, apply the triangle inequality before exchanging the sums:
\[
 \sum_i\mu_i\left|\sum_jP_{ij}(f_j-g_j)\right|
 \le\sum_{i,j}\mu_iP_{ij}|f_j-g_j|
 =\sum_j\nu_j|f_j-g_j|.
\]
For composition, positivity is termwise, the row sum is
$\sum_jP_{ij}\sum_kQ_{jk}=1$, and the new marginal is
$\sum_j\nu_jQ_{jk}=\xi_k$. Expanding a triple sum proves the action law
and associativity. The tensor equations factor into one equation for each
factor by finite Fubini. These proofs use only bounded sums and
nonnegative rational arithmetic.
\end{proof}

The existing \texttt{FiniteKernel} module concerns finite expectation
evaluation and natural-weight presentations; the matrix and completion
arguments here build on those results rather than changing their meaning
\cite{repo-finite-kernel,repo-fubini}. The general categorical setting is
related to Markov categories \cite{fritz}, with one structural restriction
explained below.

\subsection{Strong action data instead of uniform matrix approximation}

We now allow the finite probability algebras to refine indefinitely, as in
the measure presentation of the preceding section. Let $D_X,D_Y$ denote
their rational simple-function codes. Their $L^1$ completions are described
by regular names.

\begin{definition}[A presented bounded action]\label{def:action}
A bounded action $A:X\to Y$ on observables consists of a rational bound
$C\ge0$ and a PR routine $A(s,p)\in D_Y$ for $s\in D_X$. After common
refinement its finite contracts are
\begin{align}
 \norm{A(s,p)-A(t,q)}_1
 &\le C\norm{s-t}_1+2^{-p}+2^{-q},\label{eq:action-stability}\\
 \norm{A(as+bt,p)-aA(s,q)-bA(t,r)}_1
 &\le2^{-p}+|a|2^{-q}+|b|2^{-r},\label{eq:action-linearity}\\
 \norm{A(s,p)}_1&\le C\norm s_1+2^{-p}.\label{eq:action-bound}
\end{align}
They hold for all rational $a,b$, simple $s,t$, and natural precisions.
For a probability-preserving action require additionally $C=1$ and
\begin{align}
 \norm{(A(s,p))_-}_1&\le2^{-p}\quad(s\ge0),\label{eq:action-positive}\\
 \norm{A(1,p)-1}_1&\le2^{-p},\qquad
 |I_Y(A(s,p))-I_X(s)|\le2^{-p}.\label{eq:action-preserve}
\end{align}
Here $u_-=(|u|-u)/2$ is defined on simple functions and extends
nonexpansively to $L^1$. Each displayed inequality is a numerical
all-precision contract for a fixed routine.
\end{definition}

The definition does not assume a map on an already constructed uncountable
space. It gives its finite action and the precise relations needed to
construct one. For actual implementations the contracts must be proved;
placing them in a Lean structure is not a proof that an implementation
inhabits that structure.

\begin{theorem}[Extension and composition of strong actions]\label{thm:action-extension}
A presented bounded action extends uniquely to a bounded linear map on
regular $L^1$ names. Its bound is $C$. Probability-preserving actions extend
to positive unital integral-preserving contractions. Identity and
composition have PR action routines, and equality on dense simple
functions implies equality on all names.
\end{theorem}
\begin{proof}
Choose a natural $\ell$ with $C\le2^\ell$; one can take the bit length of
$\lceil\max(1,C)\rceil$. If $f$ has simple approximants $s_n$ with
$\norm{s_n-f}_1\le2^{-n}$, define
\[
 b_p=A(s_{p+\ell+2},p+2).
\]
For $p,q$ the stability contract gives
\[
 \norm{b_p-b_q}_1
 \le C(2^{-p-\ell-2}+2^{-q-\ell-2})
      +2^{-p-2}+2^{-q-2}
 \le2^{-p-1}+2^{-q-1}.
\]
Thus $b$ is a regular name. Comparing two equivalent input names with
\eqref{eq:action-stability} shows the output is independent of the name.
Taking arbitrarily precise finite comparisons gives
$\norm{Af-Ag}_1\le C\norm{f-g}_1$. The linearity contract gives exact
linearity by the same process. The negative-part, unit, and integral
contracts pass to the limit because these operations are nonexpansive in
$L^1$.

For actions $A:X\to Y$ and $B:Y\to Z$, an evaluator for $BA$ is
\[
 (B\circ A)(s,p)=B(A(s,p+\ell_B+2),p+2),
 \qquad C_B\le2^{\ell_B}.
\]
Its distance from the intended composite is at most
$C_B2^{-p-\ell_B-2}+2^{-p-2}\le2^{-p-1}$.
Triangle inequalities with the separate action contracts prove all the
required finite laws, with one extra guard bit if one insists on an exact
$2^{-p}$ convention. The bound is $C_BC_A$. The identity routine returns
its input exactly. The two associations of three actions agree on every
simple function, hence on every name by continuity. Finally, if two maps
agree on simple functions, comparison through a simple approximant $s_n$
gives output distance at most $(C_A+C_B)2^{-n}$; this proves uniqueness.
\end{proof}

\begin{example}[Why strong action is needed]\label{ex:no-operator-norm}
On the dyadic presentation of $[0,1]$, let $E_n$ average over the cells of
the $n$th partition. For every supplied $L^1$ name $f$, the previous
section constructs a PR rate for $E_nf\to f$. But for every $n$ there is
a simple function $h_n$ equal to $+1$ and $-1$ on the two halves of every
$n$th-level cell. Then $\norm{h_n}_1=1$ and $E_nh_n=0$, so
\[
 \norm{I-E_n}_{L^1\to L^1}\ge1.
\]
No sequence of these finite conditional averages approximates the identity
in operator norm. A definition that required this convergence would
exclude the identity from the intended completion. Strong action on each
named input avoids that error and provides precisely the data needed for
composition.
\end{example}

\subsection{Categorical structure and its probability constraint}

Take presented probability algebras as objects, and a transition $X\to Y$
to mean a probability-preserving action $L^1(Y)\to L^1(X)$ modulo equality
on names. Theorem~\ref{thm:action-extension} gives a category. The word
``category'' can be eliminated from any fixed argument: objects and arrows
are presentation codes, and the associativity and identity statements are
all-precision arithmetic equalities. There is no need for a universe of
all spaces inside iPRA.

For products use the rectangular finite algebra with masses $\mu_i\nu_j$.
The finite tensor matrices in Theorem~\ref{thm:finite-markov} are
contractions on this algebra. More generally, to tensor two presented
actions, write a rectangular simple function as a finite sum
$\sum_{i,j}c_{ij}\mathbf1_{A_i}\mathbf1_{B_j}$ and apply the actions
separately to the two indicators. For a product of named functions,
\[
 \norm{fg-\widetilde f\widetilde g}_{L^1(X\times Y)}
 \le\norm{f-\widetilde f}_{L^1(X)}\norm g_{L^1(Y)}
    +\norm{\widetilde f}_{L^1(X)}\norm{g-\widetilde g}_{L^1(Y)}.
\]
Finite Fubini proves this directly. It provides the PR precision schedule
for the finite tensor sum. Positivity, mass preservation, and the pointwise
inequality $|Tu|\le T|u|$ give contraction on rectangular simple
functions explicitly: put $a_i=T\mathbf1_{A_i}$ and
$b_j=S\mathbf1_{B_j}$. These functions are nonnegative, so
$|\sum c_{ij}a_ib_j|\le\sum|c_{ij}|a_ib_j$; integrating gives
$\sum|c_{ij}|\mu(A_i)\nu(B_j)$, the norm of the original atom table.
For a finite evaluator, clip indicator-image approximants to $[0,1]$.
If each has $L^1$ error at most $\delta$, each product has error at most
$2\delta$. The choice
$\delta\le2^{-p}/(2\max(1,\sum_{ij}|c_{ij}|))$ suffices for the whole
table, and only one image per source atom needs to be cached.
The tensor therefore extends to the product $L^1$ completion. Common
refinement and linearity prove independence of the chosen rectangular
decomposition. The tensor unit is the one-point probability space;
symmetry and associativity are the explicit index permutations.

This category has designated marginal measures. Copying $x\mapsto(x,x)$
generally sends $\mu$ to a diagonal law, not to $\mu\otimes\mu$.
For a fair two-point space the diagonal has mass $1$ under the copied law
and mass $1/2$ under the product law. Thus copying is not in general an
arrow to this tensor product. The full category of Markov kernels studied
in \cite{fritz} does carry copy maps; the probability-preserving category
here has the stated restricted structure. Confusing these two categories
would give incorrect independence arguments. We assert only the structure
just proved.

\subsection{Kernel perturbations and conditioning}

\begin{proposition}[A tail-sensitive kernel bound]\label{prop:kernel-tail}
Suppose $P,Q$ have the same marginals $\mu,\nu$, and set
\[
 d_\mu(P,Q)=\sum_i\mu_i\sum_j|P_{ij}-Q_{ij}|.
\]
For bounded $f$, $|f|\le K$, one has
$\norm{T_Pf-T_Qf}_{1,\mu}\le Kd_\mu(P,Q)$.
For arbitrary integrable $f$ and any $K>0$,
\begin{equation}\label{eq:kernel-tail}
 \norm{T_Pf-T_Qf}_{1,\mu}
 \le2t_K(f)+Kd_\mu(P,Q).
\end{equation}
\end{proposition}
\begin{proof}
For bounded $f$, expand the difference of the two rows and bound
$|f_j|$ by $K$ before summing. For general $f$, insert $c_K(f)$ and use
the two $L^1$ contraction bounds on $f-c_K(f)$. Each of these costs
$t_K(f)$; the remaining bounded difference costs $Kd_\mu(P,Q)$.
Completion is valid because clipping and all the operators are continuous
in $L^1$. Hence a supplied tail modulus and a kernel-approximation rate
give a strong-action rate even when no useful uniform bound on $f$ exists.
\end{proof}

\begin{proposition}[Conditioning with a stated mass margin]\label{prop:bayes}
Let $a_j,b_j\ge0$ have masses $A=\sum_ja_j$, $B=\sum_jb_j$ with
$A,B\ge\tau>0$. Put $\delta=\sum_j|a_j-b_j|$. Then
\begin{equation}\label{eq:bayes}
 \sum_j\left|\frac{a_j}{A}-\frac{b_j}{B}\right|
 \le\frac{2\delta}{\tau}.
\end{equation}
Thus bounded posterior expectations have a PR modulus once the positive
mass margin is part of the input.
\end{proposition}
\begin{proof}
Add and subtract $b_j/A$, sum, and use $\sum_jb_j=B$:
\[
 \sum_j\left|\frac{a_j}{A}-\frac{b_j}{B}\right|
 \le\frac{\delta}{A}+B\left|\frac1A-\frac1B\right|
 =\frac{\delta+|A-B|}{A}\le\frac{2\delta}{\tau}.
\]
The inequality $|A-B|\le\delta$ is finite weighted error. Clearing the
positive denominators gives a quantifier-free arithmetic proof.
\end{proof}

The margin is necessary for a uniform pointwise posterior. The unnormalized
rows $(\eps,0)$ and $(0,\eps)$ are distance $2\eps$ apart, while their
normalized laws are distance $2$ apart. Conditional expectations as $L^1$
classes instead weight each posterior by its cell mass; the cancellation
of that denominator is exactly what makes the tower construction stable.
For named real masses we do not decide whether an arbitrary mass is zero.
Finite rational masses can be tested, and continuous posterior queries
must carry a positive margin. General disintegration is an additional
construction, not a consequence of the finite tower identity alone.

\section{Continuous formulas with integration binders}\label{sec:logic}

The earlier compact-model packets already give quantitative meanings to
infimum and supremum formulas \cite{repo-abstract,repo-continuous-two}.
We extend their interface with integration. Continuous integration logic
is an established subject \cite{bagheri-pourmahdian}; the result here is a
finite-evaluation theorem with explicit hypotheses, rather than a proof of
its unrestricted compactness or ultraproduct theorems.

\subsection{Compact probability sorts and a constructive quadrature}

A \emph{compact probability sort} consists of a presented metric space
with a rational diameter bound, PR finite $h$-nets for rational $h>0$,
and a presented probability integral. Each named distance-to-center
function is supplied with compatible $L^1$ observations. The integral is
positive, takes $1$ to $1$, and has the contraction property of the measure
chapter. Compatibility is substantive: arbitrary compact and measure
presentations on unrelated carriers do not constitute one sort.

Here is an explicit sufficient construction of quadrature from this data.
It avoids having to decide the boundaries of nearest-center cells.

\begin{lemma}[Tent quadrature with rational masses]\label{lem:tent-quadrature}
Let $q_1,\ldots,q_m$ be an $h/2$-net, $m\ge1$, and set
\[
 t_i(x)=(h-d(x,q_i))_+,\quad S(x)=\sum_it_i(x),\quad
 \phi_i(x)=\frac{t_i(x)}{S(x)}.
\]
Then $S\ge h/2$, $\phi_i\ge0$, $\sum_i\phi_i=1$, and
$\phi_i(x)>0$ implies $d(x,q_i)<h$.
Let $m_i=I(\phi_i)$. If rational $a_i$ satisfy $|a_i-m_i|\le\eta$
and $m\eta\le1/2$, define
\[
 u_i=\max(0,a_i),\qquad U=\sum_iu_i,\qquad w_i=u_i/U.
\]
Then $U\ge1/2$, the $w_i$ form a rational probability vector, and
$\sum_i|w_i-m_i|\le2m\eta$.
For every $L$-Lipschitz $g$ with $|g|\le B$,
\begin{equation}\label{eq:quadrature}
 \left|I(g)-\sum_iw_i g(q_i)\right|
 \le Lh+2Bm\eta.
\end{equation}
All numerical choices are PR from the net, the integral names, and the
requested error.
\end{lemma}
\begin{proof}
Some center lies at distance at most $h/2$ from every $x$, so one tent is
at least $h/2$. This proves the positive denominator bound and the
partition identities. The reciprocal is evaluated using the globally
defined guarded function $r(s)=1/\max(h/2,s)$. It is bounded by $2/h$
and is $4/h^2$-Lipschitz. Approximate tents are clipped to $[0,h]$
before multiplication. If their $L^1$ errors are $e_j$, the resulting
error in $\phi_i$ is at most $(2/h)e_i+(4/h)\sum_je_j$. Thus intermediate
approximations need not themselves satisfy the exact lower bound on $S$.
The bounded product rule gives named $L^1$ functions $\phi_i$.

Positivity and unit mass imply $m_i\ge0$ and $\sum_im_i=1$.
Clipping $a_i$ at zero does not increase its distance from $m_i$, hence
$D=\sum_i|u_i-m_i|\le m\eta$. Thus $|U-1|\le D$ and $U\ge1/2$.
Moreover
\[
 \sum_i|u_i/U-m_i|
 \le\sum_i|u_i/U-u_i|+D
 =|1-U|+D\le2D.
\]
If $g$ is initially only a bounded Lipschitz point-evaluator, the tent
interpolants $g_h=\sum_i g(q_i)\phi_i$ first construct its $L^1$ name.
Pointwise $|g-g_h|\le Lh$; therefore any two such interpolants differ
in $L^1$ by at most the sum of their $Lh$ bounds, and a dyadic schedule
gives a regular name. This proves the needed compatibility between point
evaluation and integration rather than assuming it for every new atom.
For exact masses, use the partition identity inside the integral:
\[
 \left|I(g)-\sum_im_i g(q_i)\right|
 \le I\left(\sum_i\phi_i(x)|g(x)-g(q_i)|\right)\le Lh.
\]
The last inequality uses the support condition of each tent. Perturbing
the masses costs at most $B\sum_i|w_i-m_i|$, proving
\eqref{eq:quadrature}. Evaluating the finitely many mass names with
$\eta\le\min(1/(2m),\eps/(4Bm))$ and taking $h\le\eps/(2L)$ gives
a total error at most $\eps$, with the evident separate cases $B=0$ or
$L=0$. These are positive rational divisions and bounded loops.
\end{proof}

The tent construction itself is already present in the repository's
function-space work. The new use is to convert its support and denominator
certificates into probability quadratures with normalized rational masses.
No exact decision of a named real inequality is inserted into the algorithm.

\subsection{Syntax, moduli, and recursive evaluation}

Fix a finite many-sorted signature. Each atomic predicate is a named
bounded real-valued function of its arguments, with rational coordinate
Lipschitz bounds $L_{R,j}$ and a jointly PR evaluator on dense arguments.
Function symbols can be admitted by first computing their term moduli and
expanding them into the atomic evaluator; their evaluation error is charged
there. A formula is a finite expression generated by:
\begin{enumerate}
\item rational constants and atomic predicates;
\item rational, bounded Lipschitz connectives, including complement,
      minimum, maximum, affine averaging, and products on $[0,1]$;
\item $\inf_x\varphi$, $\sup_x\varphi$ over a specified compact sort;
\item $\mathsf E_x\varphi$, integrating over that sort's fixed probability
      measure, independent of the remaining free parameters.
\end{enumerate}
The bound-variable positions and sorts are part of the syntax. We do not
permit an unannounced change of the measure with the environment. A
parameter-dependent kernel binder requires the additional kernel-variation
contract from Section~\ref{sec:operators}.

Range bounds $B_\varphi$ and Lipschitz constants for each free variable are
computed structurally. For a connective with
$|C(a)-C(b)|\le\sum_jc_j|a_j-b_j|$, use the corresponding nonnegative
sum of the input constants. For minimum and maximum one can use the
maximum of coordinate constants. A binder drops the bound coordinate
and leaves the other constants unchanged. Integration has mass one;
infimum and supremum preserve a uniform parameter bound. The connective
range certificate supplies the next $B_\varphi$.

For a requested precision the finite evaluator replaces bound sorts by
chosen finite nets, infimum and supremum by finite minimum and maximum,
and expectation by the chosen rational weighted sum. Before using a
connective whose Lipschitz law is certified only on specified range
intervals, clip each approximate argument to its certified interval.
Clip each output to its proved range as well. Such clipping cannot
increase the error to the true value; it ensures, for example, that
the product rule on $[0,1]$ is never applied to an out-of-range sample.
The evaluator retains the syntactic nesting order. All tables are finite and all loops have their
explicit lengths. A fixed formula has a PR evaluator; a numerical syntax
code can be interpreted by bounded syntax recursion with its program
clocks, but this does not provide a universal evaluator for arbitrary PR
symbols.

\begin{theorem}[Mixed quantifier evaluation]\label{thm:mixed-evaluation}
With the preceding data, every fixed well-sorted formula defines a named
continuous function of its free variables. Suppose its atomic tables have
uniform error budgets, and each integration node uses a quadrature with
radius $h$ and mass error at most $\tau$. A sound numerical error recursion
is as follows:
\begin{align*}
 e_R&=\text{the supplied atomic evaluation error},\\
 e_{C(\varphi_1,\ldots,\varphi_s)}&=\sum_jc_je_{\varphi_j},\\
 e_{\min(\varphi,\psi)},e_{\max(\varphi,\psi)}
     &\le\max(e_\varphi,e_\psi),\\
 e_{\inf_x\varphi},e_{\sup_x\varphi}
     &\le e_\varphi+L_{\varphi,x}h,\\
 e_{\mathsf E_x\varphi}
     &\le e_\varphi+L_{\varphi,x}h+B_\varphi\tau.
\end{align*}
Here $h$ in an extremum rule is a covering radius. The output error is at
most the budget computed at the root. A PR precision schedule makes that
budget at most $2^{-p}$.
\end{theorem}
\begin{proof}
Induct on the fixed syntax. Constants and atoms are their supplied names.
The connective step is exactly its Lipschitz inequality, using the
inductively obtained range bounds. It also proves the required free-variable
moduli, so the integrands in the next steps remain uniformly continuous.

For a supremum of a function $g$ with modulus $L$, let
$M=\max_i g(q_i)$. Every $x$ is within $h$ of some $q_i$, and hence
$g(x)\le M+Lh$; conversely every $g(q_i)$ is a value of $g$. Thus the
finite maximum approximates the supremum within $Lh$. Perturbing all
table values by at most $e$ changes the maximum by at most $e$.
The infimum proof reverses the inequalities. This argument can also
\emph{construct} the supremum name without assuming its existence:
two finite nets with errors $(h,e)$ and $(h',e')$ give maxima differing by
at most $L(h+h')+e+e'$. A fixed dyadic schedule therefore yields a regular
name. Its upper-bound and approximation properties follow by comparison
with finer nets. No exact maximizer is selected.

For expectation, first replace the continuous integral by its quadrature,
costing $Lh+B\tau$ by Lemma~\ref{lem:tent-quadrature}. Replacing the
exact values at the nodes by the finite table costs at most $e$ because
the rational weights are nonnegative and normalized. This proves the
integration rule and its parameter modulus.

There are finitely many syntax nodes. The root bound is a finite
nonnegative expression in the local errors and mesh sizes. Positive
rational overestimates of its coefficients let us choose dyadic errors
small enough by bounded binary-length calculations. Each prescribed net
and mass sampler has its stated PR clock; composing them yields the
overall PR routine. For instance, make every local contribution at most
$2^{-p}/(1+s)$ after accounting for all its root influence coefficients,
where $s$ is the number of nodes. The computed bound then tends to zero
with an explicit rate. The same estimates comparing two runs give the
regularity of the resulting value name and complete the induction.
\end{proof}

If at an expectation node the table error varies with the cell, the more
accurate rule is
$\sum_iw_i e_i+Lh+B\tau$. Its soundness is
Lemma~\ref{lem:weighted}. Shared formula tables can be cached, with
their incoming influences accumulated by
Proposition~\ref{prop:adjoint}. This is where probability information
improves the numerical workload beyond a uniform worst-case schedule.

\subsection{Mixed norms and the order of binders}

The uniform rule is sufficient for every formula in the theorem. A
compiler can safely accept weaker hypotheses in special contexts.
If $\varphi(x,y)$ is followed by $\mathsf E_x$ and then $\sup_y$, an
appropriate error certificate is
\[
 \sup_y\int |\varphi(x,y)-\widetilde\varphi(x,y)|\,d\mu(x)\le e.
\]
If the order is $\sup_y$ first and $\mathsf E_x$ second, a sufficient
certificate is
\[
 \int\sup_y|\varphi(x,y)-\widetilde\varphi(x,y)|\,d\mu(x)\le e.
\]
The second is generally stronger. Their sufficiency follows from the
elementary inequalities for integration and supremum; they cannot be
interchanged merely because both contain an integral and a supremum.
Similarly, small product-space $L^1$ error by itself does not control
either supremum without a further uniformity certificate.

\begin{example}[Two genuinely different formulas]\label{ex:binder-order}
For $x,y\in[0,1]$ with Lebesgue probability measure,
\[
 \sup_y\E_x|x-y|=\frac12,
 \qquad
 \E_x\sup_y|x-y|=\frac34.
\]
Indeed $\E_x|x-y|=y^2/2+(1-y)^2/2$, whose maximum on $[0,1]$ is
$1/2$. On the other hand $\sup_y|x-y|=\max(x,1-x)$, with integral
$3/4$. On the $m$ equally weighted midpoint nodes, for even $m$, the
two finite values are respectively
\[
 \frac12-\frac1{2m},\qquad\frac34-\frac1{2m}.
\]
For the first, the extreme midpoint maximizes a convex average of
absolute values, and direct summation gives $(m-1)/(2m)$. For the second,
the maximum is $\max(x,1-x)-1/(2m)$ on each midpoint; the symmetric
midpoint sum of $|x-1/2|$ is $1/4$. Thus the finite evaluator preserves
the nonzero gap exactly. The supplied reference program computes both
syntax trees independently.
\end{example}

Two expectation binders over independent product measures can be
exchanged, with proof by finite Fubini followed by the named limit.
Their exchange is a proved rule, not a general permission to reorder
binders. The example shows why the syntax must retain the difference.

\subsection{Model-theoretic boundaries}

The metric space of all events of an atomless probability algebra,
with distance $\mu(A\mathbin\triangle B)$, need not be totally bounded.
Independent fair-coordinate events have pairwise distance $1/2$.
Thus the compact-net theorem does not evaluate unrestricted quantifiers
over all events. It evaluates compact point sorts, supplied compact
families of observables, or explicitly finite event presentations.
The model theory of probability algebras has substantially richer
structure \cite{berenstein-henson}; its full quantifier elimination and
model existence statements are not consequences of this finite theorem.

A computed value interval can certify a strict inequality separated by
a rational margin. It cannot in general decide exact zero truth.
Similarly, values of every finite test do not by themselves select one
coherent model of an arbitrary theory. For that task one needs compatible
named structure data or an additional constructive extension procedure.
These are mathematical boundaries of the embedding, not defects that can
be removed by replacing words with certificates.

\section{Holomorphic parameter integrals from integrable names}\label{sec:complex-integrals}

The measure construction also supports analysis of continuous parameters.
The relevant principle is that a proved bounded linear map commutes with a
series having a supplied norm-tail schedule. This is stronger than simply
approximating one finite integral, and it does not require a pointwise
decision of convergence on a null set.

\subsection{Series in a presented normed space}

Let $B$ be the completion of a PR-presented normed rational vector space;
for complex scalars use rational pairs and a compatible submultiplicative
norm. Suppose $a_n\in B$ have jointly PR names and rational data $M,R>0$
satisfy
\begin{equation}\label{eq:coefficient-majorant}
 \norm{a_n}_B R^n\le M\qquad(n\in\N).
\end{equation}
For $|z|\le r<R$, put $q=r/R$. Define finite sums
$F_N(z)=\sum_{n<N}a_nz^n$.

\begin{lemma}[Norm-tail schedule]\label{lem:banach-series}
For $N\le L$,
\begin{equation}\label{eq:banach-tail}
 \norm{F_L(z)-F_N(z)}_B
 \le M\frac{q^N}{1-q}.
\end{equation}
Thus $F(z)=\sum_{n\ge0}a_nz^n$ has a regular name, uniformly on the
closed disk $|z|\le r$. If a finite approximation $\widetilde a_n$ has
error $\eta_n$, the extra finite-sum error is at most
$\sum_{n<N}r^n\eta_n$.
\end{lemma}
\begin{proof}
Apply the triangle inequality and the coefficient bound termwise:
$\sum_{n=N}^{L-1}\norm{a_n}|z|^n\le M\sum_{n=N}^{L-1}q^n$.
The finite geometric identity bounds this by
$Mq^N/(1-q)$. The second assertion is the same weighted triangle
inequality applied to coefficient errors.

Here is a PR clock without logarithmic search. For a rational $0\le q<1$,
let $s=\lceil1/(1-q)\rceil$. Induction gives
$q^j\le1/(1+j(1-q))$, so $q^s\le1/2$ and $q^{st}\le2^{-t}$.
Choose a natural $b$ with $M/(1-q)\le2^b$, using the binary length of
its ceiling. Then $N=s(p+b+2)$ makes the tail at most $2^{-p-2}$.
Evaluate the finite coefficients so that their weighted error is at most
$2^{-p-2}$; for example a common error
$2^{-p-2}/(1+\sum_{n<N}r^n)$ suffices. Comparison of two such runs
gives a regular name. All loops and all rational operations have explicit
PR bounds. The case $q=0$ can instead use one term exactly.
\end{proof}

For a named $z$ with $|z|\le r$, evaluate a rational approximation inside
a larger certified radius $s_0=(r+R)/2$. A precision making its error
at most $(R-r)/2$ guarantees that bound. The derivative estimate below
supplies a Lipschitz constant for the additional input error. This is a
positive radius-gap calculation, not an exact test of whether a named
complex number lies on a circle.

\subsection{Integration and differentiation}

Take $B=L^1(\mu;\C)$ with its usual complex norm, obtained by completing
complex simple functions. The complex integral is obtained by integrating
real and imaginary parts; its norm is at most the $L^1$ norm, and
$\norm{zv}_B=|z|\norm v_B$. These are the norms used in all disk
estimates below. Exact rational calculations may use
$|\cdot|_\square$ as an upper certificate, since
$|z|\le|z|_\square\le2|z|$. A conversion back to a square-norm error
therefore incurs a factor at most $2$, which must be budgeted explicitly;
the same disk constants are not silently asserted for that different
norm. Rational upper certificates or rational square-root enclosures
give PR presentations of the usual norm with these guard factors.

\begin{theorem}[Constructive holomorphic parameter integration]\label{thm:holomorphic-integral}
Under \eqref{eq:coefficient-majorant}, let $f(z)=\sum_na_nz^n$ as an
$L^1$-valued name. Then the scalar function $H(z)=I(f(z))$ satisfies
\begin{equation}\label{eq:integral-series}
 H(z)=\sum_{n\ge0}I(a_n)z^n\qquad(|z|<R).
\end{equation}
It is holomorphic, and for every requested natural derivative order $j$,
\begin{equation}\label{eq:integral-derivative}
 H^{(j)}(z)=I\left(\sum_{m\ge0}\frac{(m+j)!}{m!}
                    a_{m+j}z^m\right).
\end{equation}
The series, its derivatives, and the identities have uniform PR
finite-observation schedules on every supplied smaller radius $r<R$.
\end{theorem}
\begin{proof}
Linearity gives the identity for every finite sum. The integral contraction
bounds the difference between the integral of the limiting name and the
finite scalar sum by \eqref{eq:banach-tail}. Consequently both sides of
\eqref{eq:integral-series} are the same regular scalar name. This already
justifies the exchange of integration and the infinite sum.

To justify differentiation directly, choose $r<s_0<R$. For
$|z|\le r$ and $|z|+|h|\le s_0$, the finite polynomial identity for
the second-order remainder gives
\[
 |(z+h)^n-z^n-nz^{n-1}h|
 \le \binom n2s_0^{n-2}|h|^2\qquad(n\ge2).
\]
One proof is to expand the telescoping difference twice; the remaining
sum has $\binom n2$ terms, each a product of $n-2$ factors bounded by
$s_0$. Multiplying by the coefficient majorant and summing yields
\begin{equation}\label{eq:holomorphic-remainder}
 \norm{f(z+h)-f(z)-h\sum_{n\ge1}na_nz^{n-1}}_1
 \le\frac{M|h|^2}{R^2(1-s_0/R)^3}.
\end{equation}
The finite identity
$\sum_{n\ge2}\binom n2t^{n-2}=(1-t)^{-3}$ is read here through its
finite partial sums and geometric tails; alternatively the needed upper
bound follows from multiplying three nonnegative geometric sums and
counting triples. Thus it has a bounded-sum proof. Dividing by $|h|>0$
gives an explicit difference-quotient modulus. Applying the integral
contraction proves complex differentiability of $H$ with the stated first
derivative. To repeat the argument, for any $r<R'<R$ the derivative
coefficient sequence $b_m=(m+j)!a_{m+j}/m!$ has majorant
\[
 \norm{b_m}_1(R')^m\le
 \frac{Mj!}{R^j(1-R'/R)^{j+1}}.
\]
This follows from the nonnegative binomial-generating bound used below.
Thus the differentiated series again satisfies the hypotheses on the
smaller radius, and the same argument proves every finite order.

For an explicit tail for order $j$, put
\[
 q=r/R,\qquad\theta=(1+q)/2,\qquad\alpha=q/\theta<1.
\]
If $q>0$, then
\begin{align}
 \sum_{m\ge N}\frac{(m+j)!}{m!}\norm{a_{m+j}}_1r^m
 &\le\frac{Mj!}{R^j}
       \sum_{m\ge N}\binom{m+j}{j}q^m\notag\\
 &\le \frac{Mj!}{R^j(1-\theta)^{j+1}}
                  \frac{\alpha^N}{1-\alpha}.\label{eq:derivative-tail}
\end{align}
Indeed $\binom{m+j}{j}\theta^m\le(1-\theta)^{-j-1}$,
as seen by taking one term of the nonnegative binomial generating series.
The remaining geometric series proves the second inequality. All constants
are rational and factorials are PR, so Lemma~\ref{lem:banach-series}'s
block clock applies to $\alpha$. When $r=0$, only the term $m=0$ is
needed and $H^{(j)}(0)=j!I(a_j)$.

Finally, evaluating the finitely many derivative coefficients with errors
$\eta_{m+j}$ costs at most
$\sum_{m<N}(m+j)!/m!\,r^m\eta_{m+j}$. This supplies a direct weighted
precision schedule for \eqref{eq:integral-derivative}. No convergence
principle without a norm-tail bound has been used.
\end{proof}

This theorem reuses standard geometric/derivative arguments from the
repository's complex-series interfaces \cite{repo-complex-series}.
Its additional ingredient is the constructed $L^1$ integral and its
contraction property. In particular, the integrand may have an unbounded
spatial factor. If $g\in L^1[0,1]$ has a supplied name, put
$a_n(t)=g(t)t^n$. Bounded multiplication supplies jointly PR coefficient
names, and $\norm{a_n}_1\le\norm g_1$. Hence
\[
 H(z)=\int_0^1\frac{g(t)}{1-zt}\,dt\qquad(|z|<1)
\]
has the stated constructive analytic interpretation. To identify the
quotient with the named series, use the exact finite identity
\[
 \frac{g(t)}{1-zt}-g(t)\sum_{n<N}(zt)^n
 =\frac{g(t)(zt)^N}{1-zt}.
\]
For $|z|\le r<1$ its $L^1$ norm is at most
$\norm g_1 r^N/(1-r)$. The reciprocal is a bounded multiplier because
$|1-zt|\ge1-r$, so this also constructs the quotient's $L^1$ name.
A rational upper
bound for $\norm g_1$ suffices; a pointwise bound for $g$ is unnecessary.

\subsection{A complete rational example}

For $g=1$, the coefficient integral is $I(t^n)=1/(n+1)$, and
\begin{equation}\label{eq:log-example}
 H(z)=\int_0^1\frac{dt}{1-zt}
     =\sum_{n\ge0}\frac{z^n}{n+1}.
\end{equation}
The identity for $I(t^n)$ has an entirely finite proof. Factor
\[
 (k+1)^{n+1}-k^{n+1}
 =\sum_{j=0}^{n}(k+1)^{n-j}k^j.
\]
Each summand lies between $k^n$ and $(k+1)^n$.
Sum for $0\le k<m$ and divide by $(n+1)m^{n+1}$. The resulting
inequality places $1/(n+1)$ between the left and right Riemann sums for
$t^n$ on the $m$-mesh. Their difference is exactly $1/m$ for $n\ge1$
and zero for $n=0$. Hence they construct the asserted integral with
the PR mesh $m=2^{p+2}$. Factoring $x^n-y^n$ also gives the uniform
cell-oscillation bound $n/m$, useful when constructing the full $L^1$
name of $t^n$. The reference program checks these rational brackets
and the finite series identities.

On $|z|\le1/2$ the $N$-term scalar approximation has error at most
$2^{1-N}$. The value at the center is $H(0)=1$. If the analytic logarithm
on $|z|<1$ is defined by its usual series, then for $z\ne0$
$H(z)=-\log(1-z)/z$. Formula~\eqref{eq:log-example} is the definition
used at zero, so it introduces neither a spurious singularity nor a
branch-selection problem.

The same construction gives
$H^{(j)}(0)=j!/(j+1)$. This is an equality of named complex numbers;
for the rational coefficient values it is an exact arithmetic identity.
The code checks both series-tail and derivative-tail certificates, while
the general statements are proved above.

\section{Quantitative homological perturbation}
\label{sec:homological}

The effective nerve arguments in the existing packets obtain a small global
model when all relevant intersections are acyclic or carry supplied
contractions.  A useful extension retains a small model of each intersection
even when its homology is nontrivial.  It must transport the differential,
cycles, boundaries and filling operators together.  Merely approximating
boundary matrices does not suffice: an approximately square-zero matrix is
not a chain complex.

The algebraic tool is the classical homological perturbation lemma
\cite{crainic2004}.  ProveIt's unknot-recognition report~13 already implements
a finite version over $\mathbb F_2$, with a nilpotence bound coming from a
positive grading~\cite{repo_hpl}.  The contribution here is its integration
with quantitative presentations: signed rational formulas, an entirely
finite residual theorem, explicit precision schedules, and a bounded-operator
interpretation.  The underlying perturbation lemma is established
mathematics.  The new interface avoids repeatedly rebuilding a homological
model when its differential changes.

\subsection{Data, signs and an exact algebraic theorem}

A bounded chain complex is a finite family of vector spaces $C_j$ with maps
$d_j:C_j\to C_{j-1}$ satisfying $d_{j-1}d_j=0$.  Write $C$ for their direct
sum and $d$ for the resulting degree $-1$ map.  Initially the spaces can be
finite-dimensional rational coordinate spaces; later they may be Banach
spaces with bounded structural maps.  A smaller complex $(B,b)$ comes with
degree-zero maps $i:B\to C$, $p:C\to B$ and a degree-one map $h:C\to C$.
Our input is the complete list of identities
\begin{align}
 d^2=b^2&=0,& di&=ib,& pd&=bp,& pi&=1_B,\\
 dh+hd&=1_C-ip,& hi&=0,& ph&=0,& h^2&=0.
 \label{eq:hpl-input}
\end{align}
These data are a normalized strong deformation retraction.  Multiplication
of operators means composition, with the rightmost operator acting first.
The sign in $dh+hd=1-ip$ fixes every sign below.  In the convention
$ip=1+dh+hd$ of~\cite{crainic2004}, replace $h$ by $-h$ to obtain ours.

A perturbation is a degree $-1$ map $\delta:C\to C$ for which
\begin{equation}
 (d+\delta)^2=0,
 \quad\text{equivalently}\quad d\delta+\delta d+\delta^2=0.
 \label{eq:hpl-square-zero}
\end{equation}
It is important that this is an exact hypothesis.  Put $d'=d+\delta$.

\begin{theorem}[Signed perturbation with explicit witnesses]
\label{thm:hpl-exact}
Assume \eqref{eq:hpl-input}--\eqref{eq:hpl-square-zero}.  Suppose that
$1+h\delta$ and $1+\delta h$ have two-sided inverses $F,G$, respectively.
Define
\begin{align}
 A&=G\delta=\delta F,&
 b'&=b+pAi,\\
 i'&=i-hAi=Fi,&
 p'&=p-pAh=pG,&
 h'&=h-hAh=Fh=hG.
 \label{eq:hpl-transferred}
\end{align}
Then $(B,b')$ is a complex and
\begin{align}
 (b')^2&=0,& d'i'&=i'b',&p'd'&=b'p',&p'i'&=1_B,\\
 d'h'+h'd'&=1_C-i'p',&h'i'&=0,&p'h'&=0,&(h')^2&=0.
 \label{eq:hpl-output}
\end{align}
If $(h\delta)^L=(\delta h)^L=0$ for a supplied positive integer $L$, the
inverses and all witnesses are finite rational expressions:
\[
 F=\sum_{j=0}^{L-1}(-h\delta)^j,
 \qquad G=\sum_{j=0}^{L-1}(-\delta h)^j.
\]
\end{theorem}

\begin{proof}
The identity $(1+\delta h)\delta=\delta(1+h\delta)$ gives $G\delta=\delta F$.
From the two inverse equations,
\begin{equation}
 A=\delta-\delta hA=\delta-Ah\delta,
 \qquad F=1-hA,\qquad G=1-Ah.
 \label{eq:hpl-resolvent}
\end{equation}
Put $e=ip$, $L_0=1+\delta h$ and $R_0=1+h\delta$.  A key cancellation is
\begin{align*}
 L_0(dA+Ad)R_0
 &=d\delta+\delta d+\delta(hd+dh)\delta\\
 &=-\delta^2+\delta(1-e)\delta=-\delta e\delta.
\end{align*}
Multiplication by the two inverses therefore yields
\begin{equation}
 dA+Ad+AeA=0.
 \label{eq:hpl-maurer}
\end{equation}
This is an algebraic identity, not a convergence argument.  It proves
square-zero of the new differential directly:
\[
 (b+pAi)^2
 =b^2+p(dA+Ad+AeA)i=0.
\]
The two original chain-map identities were used to move $b$ across $p$ and
$i$ in this expansion.

For the inclusion, use $dh=1-e-hd$ and \eqref{eq:hpl-resolvent} to compute
\[
 d'(i-hAi)=ib+eAi+hdAi.
\]
Equation~\eqref{eq:hpl-maurer} changes the last term into
$-hAdi-hAeAi$, which gives $(i-hAi)(b+pAi)$.  Similarly
\[
 (p-pAh)d'=bp+pAe+pAdh
              =(b+pAi)(p-pAh).
\]
Expanding $(p-pAh)(i-hAi)$ leaves $pi=1_B$: every other term contains
$ph$, $hi$ or $h^2$.

For the homotopy identity, first absorb the terms containing a single
$\delta$ using \eqref{eq:hpl-resolvent}.  The full calculation is
\begin{align*}
 d'h'+h'd'
 &=1-e+Ah+hA-dhAh-hAhd\\
 &=1-e+eAh+hAe+h(dA+Ad)h\\
 &=1-e+eAh+hAe-hAeAh\\
 &=1-(i-hAi)(p-pAh).
\end{align*}
The three remaining side conditions follow by expansion; each summand
contains one of $h^2$, $hi$, $ph$.  Finally,
$(1+T)\sum_{j<L}(-T)^j=1-(-T)^L$, on either side.  A supplied
nilpotence exponent therefore supplies the two inverses by finite sums.
\end{proof}

The theorem gives a quotient-free formulation of homology preservation.
If $d'z=0$, then $b'p'z=0$ and
\begin{equation}
 z-i'p'z=d'h'z.
 \label{eq:hpl-cycle-transfer}
\end{equation}
If $b'w=0$, then $d'i'w=0$ and $p'i'w=w$.  If $i'w=d'v$, then
$w=b'p'v$.  These are equations on vectors and explicit witnesses; the
object language need not construct the quotient $\ker d'/\operatorname{im}d'$.

\subsection{Finite residuals without an infinite inverse}

Equip finite rational coordinate spaces with the maximum norm and matrices
with the maximum absolute row-sum norm.  This norm is rational, is computed
by bounded folds, and is submultiplicative.  The same argument applies to
any submultiplicative operator norm.  Let the supplied nonnegative rational
bounds be
\[
 \|d\|\le D,\quad \|h\|\le H,\quad \|\delta\|\le U,
 \quad\|i\|\le I,\quad\|p\|\le P,
 \quad\|h\delta\|,\|\delta h\|\le q<1.
\]
One can use $q=HU$ when that is smaller than one; computing the two products
can give a much sharper certificate.  No inverse, spectrum, or rank is
part of this input.

For $N\ge0$ put
\begin{align}
 A_N&=\delta\sum_{j=0}^{N-1}(-h\delta)^j,
 &R_N&=\delta(-h\delta)^N=(-\delta h)^N\delta,\\
 M_N&=dA_N+A_Nd+A_NipA_N.
 \label{eq:hpl-finite-data}
\end{align}
The empty sum is zero.  Define $b_N,i_N,p_N,h_N$ by the formulas
\eqref{eq:hpl-transferred}, using $A_N$ instead of $A$.

\begin{theorem}[A finite homological residual certificate]
\label{thm:hpl-finite}
Let
\begin{equation}
 K_M=\frac{U\bigl(2(1+q)D+3UIP\bigr)}{(1-q)^2}.
 \label{eq:hpl-K}
\end{equation}
For every $N\ge0$, the following identities hold exactly:
\begin{align}
 p_Ni_N&=1_B,&h_Ni_N&=0,&p_Nh_N&=0,&h_N^2&=0,\label{eq:hpl-finite-normalizations}\\
 b_N^2&=pM_Ni,\label{eq:hpl-residual-first}\\
 d'i_N-i_Nb_N&=R_Ni+hM_Ni,\\
 p_Nd'-b_Np_N&=pR_N+pM_Nh,\\
 d'h_N+h_Nd'-(1-i_Np_N)&=R_Nh+hR_N+hM_Nh.
 \label{eq:hpl-residual-identities}
\end{align}
Moreover $\|M_N\|\le K_Mq^N$, and consequently the four potentially nonzero residuals
in \eqref{eq:hpl-residual-first}--\eqref{eq:hpl-residual-identities} are bounded, in their displayed order,
by
\begin{equation}
 PI K_M q^N,\quad
 I(U+HK_M)q^N,\quad
 P(U+HK_M)q^N,\quad
 (2HU+H^2K_M)q^N.
 \label{eq:hpl-residual-bounds}
\end{equation}
All assertions concern finite arithmetic data.
\end{theorem}

\begin{proof}
The finite geometric identity gives
\[
 \delta-A_N-\delta hA_N
 =\delta-A_N-A_Nh\delta=R_N,
 \qquad \|R_N\|\le Uq^N.
\]
Repeating the expansions in Theorem~\ref{thm:hpl-exact}, retaining these
two remainder terms, gives every identity in
\eqref{eq:hpl-finite-normalizations}--\eqref{eq:hpl-residual-identities}.  In particular the new differential's
square is a computed defect; it is not silently assumed to vanish.

Here is a direct finite bound on that defect's common source.  With
$e=ip$, $L_0=1+\delta h$, $R_0=1+h\delta$, multiplication gives
\begin{align*}
 L_0M_NR_0={}&-L_0dR_N-R_NdR_0
              -\delta eR_N-R_Ne\delta+R_NeR_N.
\end{align*}
For any matrix $X$, the triangle inequality gives
\[
 (1-q)\|X\|\le\|L_0X\|,
 \qquad (1-q)\|X\|\le\|XR_0\|.
\]
Thus $(1-q)^2\|M_N\|\le\|L_0M_NR_0\|$, without constructing an
inverse.  Use $\|e\|\le IP$, $\|L_0\|,\|R_0\|\le1+q$ and
$q^N\le1$ to obtain
\[
 (1-q)^2\|M_N\|
 \le \bigl(2(1+q)D+3UIP\bigr)Uq^N.
\]
Dividing by the explicitly positive rational $(1-q)^2$ proves the claim.
Submultiplicativity applied to the exact residual identities gives the four
remaining bounds.
\end{proof}

The constants in this finite theorem are deliberately conservative.  Its
advantage is that the arithmetic proof never quantifies over an infinite
operator or uses completeness.  A finite truncation generally has
$b_N^2\ne0$; until the residual is exactly zero it is an approximation to a
complex, not a new exact complex.

\subsection{Primitive-recursive clocks and the continuous interpretation}

\begin{lemma}[A rational geometric clock]
\label{lem:hpl-clock}
For $0\le q<1$ and rational $K\ge1$, let
$s=\lceil(1-q)^{-1}\rceil$.  Choose an integer $c$ with $K\le2^c$,
for example the bit length of the numerator in a positive rational code
for $K$.  Then
\[
 N=s(k+c)\quad\Longrightarrow\quad Kq^N\le2^{-k}.
\]
Every operation in this schedule is primitive recursive on rational codes.
\end{lemma}

\begin{proof}
We have $q\le1-1/s$.  The polynomial inequality
$(1-t)^n(1+nt)\le1$ for $0\le t\le1$ follows by induction: multiplying
the next factor uses
$(1-t)(1+(n+1)t)=1+nt-(n+1)t^2\le1+nt$.
At $t=1/s$ and $n=s$ this gives $q^s\le1/2$, including $s=1$.
Taking $k+c$ blocks and multiplying by $K\le2^c$ gives the result.
After denominator clearing these are quantifier-free rational inequalities
and induction on natural-number counters.
\end{proof}

Take $K$ to be at least one and at least the four constants in
\eqref{eq:hpl-residual-bounds}.  This supplies a single explicit clock for
all homological equations.  The original coordinate dimensions, the
iteration bound $N$, and the rational operations determine every loop.
There is no search for a stage at which an unquantified limit is sufficiently
close.

\begin{theorem}[Bounded-operator completion and stability]
\label{thm:hpl-banach}
Suppose the spaces in \eqref{eq:hpl-input} are Banach spaces, all displayed
maps are bounded, and the same rational norm bounds and $q<1$ are supplied.
Then the series defining $F,G$ converge, Theorem~\ref{thm:hpl-exact} applies,
and
\begin{align}
 \|A-A_N\|&\le\frac{Uq^N}{1-q},\\
 \|i'-i_N\|&\le\frac{Iq^{N+1}}{1-q},&
 \|p'-p_N\|&\le\frac{Pq^{N+1}}{1-q},\\
 \|h'-h_N\|&\le\frac{Hq^{N+1}}{1-q},&
 \|b'-b_N\|&\le\frac{PIUq^N}{1-q}.
 \label{eq:hpl-tail}
\end{align}
Also $\|i'\|\le I/(1-q)$, $\|p'\|\le P/(1-q)$ and
$\|h'\|\le H/(1-q)$.

For two perturbations $\delta_1,\delta_2$ of the same original contraction,
suppose both have the product bound $q$ and
$\|\delta_1-\delta_2\|\le\eta$.  Then
\begin{equation}
 \|A_1-A_2\|\le\frac{\eta}{(1-q)^2}.
 \label{eq:hpl-stability}
\end{equation}
The corresponding differences of $b',i',p',h'$ are bounded by this
quantity multiplied by $PI,HI,PH,H^2$, respectively.
\end{theorem}

\begin{proof}
For any finite $M$ the tails of the geometric sums have norm at most
$q^N\sum_{j<M}q^j\le q^N/(1-q)$.  Thus the explicit clock above makes
the sums Cauchy.  Completeness supplies their limits in the bounded-operator
space.  The finite inverse defects are $(-h\delta)^N$ and
$(-\delta h)^N$, tending to zero; continuity of multiplication gives the
two-sided inverse equations.  This proves existence with a rate, rather
than appealing to a qualitative inverse theorem.

The first tail bound follows on multiplying by $\delta$.  The identity
$1-hA_N=\sum_{j=0}^{N}(-h\delta)^j$ shows why the inclusion and homotopy
bounds have exponent $N+1$; the projection has the analogous identity.
Multiplying by $p,i,h$ gives \eqref{eq:hpl-tail}.  Summing the geometric
majorant gives the norm bounds.

Write $G_j=(1+\delta_jh)^{-1}$ and $F_j=(1+h\delta_j)^{-1}$.
The resolvent identity is
\[
 A_1-A_2=G_1(\delta_1-\delta_2)F_2.
\]
Indeed $G_1-G_2=-G_1(\delta_1-\delta_2)hG_2$; substitute this in
$G_1\delta_1-G_2\delta_2$ and use $1-hA_2=F_2$.
The two inverse norm bounds are $1/(1-q)$, proving
\eqref{eq:hpl-stability}.  The remaining assertions follow from
\eqref{eq:hpl-transferred}.
\end{proof}

This Banach statement is the semantic interpretation.  Its iPRA content is
the family of finite tail and residual inequalities, represented by rational
codes and proved with the stated clocks.  To use infinite-dimensional
operators computationally, one additionally supplies a rational operator
presentation with primitive-recursive composition, evaluation and certified
norm bounds, or a fixed algorithm with proofs of these properties.
An arbitrary bounded operator is not automatically such an input.
Unbounded differentials require separate invariant-domain and graph-norm
arguments and are outside this theorem.

\subsection{Non-acyclic overlaps and quantitative fillings}

\begin{corollary}[A global model from supplied local models]
\label{cor:hpl-bicomplex}
Let $C_{a,j}$, $0\le a\le r$, be a finite-column double complex.  Its
vertical differential $d$ lowers $j$, its horizontal differential $\delta$
lowers $a$, and
$d^2=\delta^2=d\delta+\delta d=0$.  Suppose each column carries data
\eqref{eq:hpl-input} onto a supplied smaller complex, with $i,p,h$
preserving the column number.  Then their direct sum retracts the total
complex onto the smaller total space with differential
\begin{equation}
 b'=b+\sum_{j=0}^{r}(-1)^j p\delta(h\delta)^ji.
 \label{eq:hpl-global}
\end{equation}
The retraction maps and all cycle and boundary witnesses are finite
primitive-recursive expressions.  No local acyclicity assumption is needed.
\end{corollary}

\begin{proof}
Both $h\delta$ and $\delta h$ lower the column number.  Their $(r+1)$st
powers therefore vanish: on every basis block, no nonzero target column
remains.  Apply the nilpotent case of Theorem~\ref{thm:hpl-exact} with
$L=r+1$.  The displayed formula is $b+pAi$ with that finite inverse.
All chain identities and the explicit equations
\eqref{eq:hpl-cycle-transfer} follow from the same theorem.  A signed
totalization of a geometric double complex must supply the stated
anticommutation identity; it is not inferred from approximate matrices.
\end{proof}

For a finite cover this allows each intersection to retain its own explicit
homological model.  The terms with $j>0$ account for successive excursions
through eliminated pieces.  When intersections have nontrivial homology,
the result is generally larger than the ordinary nerve.  Thus the earlier
acyclic-overlap restriction is replaced by computable local models, with
the corrected global differential recorded explicitly.  Constructing those
local models is a remaining geometric task for each application.

Here is a precise finite-cover instance, including its comparison with the
original space.  Let a finite ordered simplicial complex $K$ be the union
of subcomplexes $K_1,\ldots,K_m$.  In bidegree $(a,j)$ take the direct sum
of rational simplicial $j$-chains on $K_{i_0}\cap\cdots\cap K_{i_a}$,
over strictly increasing index tuples.  The horizontal map deletes a cover
index with the alternating sign.  The vertical map is $(-1)^a$ times the
simplicial boundary.  The two maps anticommute because deleting a cover
index changes the parity of $a$, while restriction commutes with the
ordinary simplicial boundary.

For each simplex $\sigma$ of $K$, the nonempty set of cover indices
containing it has a least member $u(\sigma)$.  Its horizontal row is the
chain complex of a full simplex on those indices.  Coning an index tuple
to $u(\sigma)$ contracts this row onto one copy of $\sigma$: insert
$u(\sigma)$ into the ordered tuple with the insertion sign, and return
zero if it was already present.  In the boundary of this cone, the term
deleting the inserted vertex is the original tuple; every other deletion
cancels its counterpart obtained by coning after deletion.  In degree zero
the remaining term is the chosen vertex.  This proves
$\delta h_0+h_0\delta=1-i_0p_0$.  Also $h_0^2=h_0i_0=p_0h_0=0$,
since insertion is performed at most once and raises the horizontal degree.
Here $p_0$ forgets the singleton cover index and is zero in higher
horizontal degree, and $i_0$ chooses $u(\sigma)$.

Apply Theorem~\ref{thm:hpl-exact} to this horizontal contraction, with the
vertical boundary as perturbation.  Each product of that boundary with
$h_0$ lowers the simplicial degree, so its $(\dim K+1)$st power is zero
for a nonempty complex.  Every positive correction term in the transferred
small differential raises the horizontal degree before applying $p_0$,
and is therefore killed.  The surviving term $p_0di_0$ is exactly the
simplicial boundary of $K$.  Thus the total double complex contracts onto
$C_*(K;\Q)$ by explicit finite formulas.  Applying
Corollary~\ref{cor:hpl-bicomplex} in the other direction, using the supplied
intersection models, gives an explicit homology comparison between those
updated small models and $K$.  The empty complex is handled by the zero
complex.  No appeal to a spectral sequence or an unconstructed comparison
map is needed.  This also clarifies the limitation: the model may be large
when the intersections have large homology, and no uniform speed advantage
over direct elimination is asserted.

The continuous version also gives useful estimates on actual fillings.
If $d'z=0$ and a vector $s$ in the small model satisfies $b's=p'z$, then
\[
 v=h'z+i's,\qquad d'v=z,
 \qquad
 \|v\|\le\frac{H\|z\|+I\|s\|}{1-q}.
\]
The equation follows by adding the two chain identities.  In particular,
when $B=0$, every cycle has the canonical filling $h'z$ with norm at most
$H\|z\|/(1-q)$.  Conversely, a supplied linear functional $\lambda$ on
$B$ with $\lambda b'=0$ and $\lambda p'z\ne0$ proves that $z$ cannot be a
boundary.  A positive rational lower bound for $|\lambda p'z|$ makes this
obstruction robust under certified approximation.  These are explicit
filling or obstruction certificates, without an undecidable test of equality
of arbitrary real coefficients.

\subsection{A continuous Schur-complement example and a necessary boundary}

Take a two-term complex $C_1=C_0=\Q^2$ and
\[
 d=\begin{pmatrix}1&0\\0&0\end{pmatrix},\qquad
 B_1=B_0=\Q,\quad b=0.
\]
Let $i$ include the second coordinate, $p$ project to it, and let
$h:C_0\to C_1$ be the matrix $\operatorname{diag}(1,0)$, with $h=0$ on
$C_1$.  These satisfy all the normalized contraction identities on the
graded direct sum.  Perturb the only differential by
\[
 \delta=\begin{pmatrix}a&\beta\\\gamma&\tau\end{pmatrix}.
\]
There is no next differential, so square-zero of $d+\delta$ is automatic.
Whenever the supplied margin makes $1+a$ invertible, the transferred
differential is
\begin{equation}
 b'=\tau-\frac{\gamma\beta}{1+a}.
 \label{eq:hpl-schur}
\end{equation}
Indeed $i'$ on degree one is the column
$(-\beta/(1+a),1)^{\mathsf T}$, $p'$ on degree zero is the row
$(-\gamma/(1+a),1)$, and the surviving entry of $h'$ is $1/(1+a)$.
Multiplication by $d+\delta$ gives \eqref{eq:hpl-schur} directly.

For $(a,\beta,\gamma,\tau)=(1/4,1/8,1/5,1/50)$, the small differential
is zero and the product norms are bounded by $q=3/8$.  Changing only
$\tau$ to $1/25$ changes $b'$ to $1/50$.  Thus a small differential
perturbation can destroy homology on the retained model, even though all
perturbation formulas remain well-conditioned.  The theorem preserves
equivalence to the \emph{updated} small complex, not the original Betti
numbers.  Stable Betti numbers require extra information, such as an
appropriate rank margin or vanishing of the transferred perturbation.

The exact-rational executable \texttt{code/homological\_examples.py}
checks this example against an independently computed rational inverse,
verifies all contraction identities, and checks the finite residual bounds.
It also treats $a=0$, where both product squares vanish and the two-term
finite series is exact even without a norm-smallness premise.

\subsection{A precise Lean and object-language development path}

The smallest semantic interface is a structure of rational matrices with
fields $d,b,i,p,h$, shape and degree data, and proofs of
\eqref{eq:hpl-input}.  A second structure adds $\delta$ and the proof of
\eqref{eq:hpl-square-zero}.  Define \texttt{neumann} by primitive recursion
on $N$, and define $A_N,R_N,M_N$ by \eqref{eq:hpl-finite-data}.
The main theorem dependencies are:
\begin{enumerate}
\item \texttt{geometric\_defect}: both finite inverse defects equal the
appropriate $N$th power; prove by induction on $N$.
\item \texttt{residual\_identities}: the seven groups of identities in
\eqref{eq:hpl-finite-normalizations}--\eqref{eq:hpl-residual-identities}; prove by distributivity and the supplied
chain and normalization equations.  This is noncommutative algebra, so a
commutative \texttt{ring} call on whole matrices is inappropriate.
\item \texttt{rowNorm\_mul\_le}: the rational maximum-row-sum norm is
submultiplicative; prove finite triangle bounds, then the bounded maximum
property.  Use it to prove \texttt{finite\_residual\_bound} exactly as in
Theorem~\ref{thm:hpl-finite}.
\item \texttt{block\_clock}: prove Lemma~\ref{lem:hpl-clock} over rational
codes after denominator clearing.  The invariant is the displayed
polynomial inequality, not a real logarithm.
\item \texttt{filtered\_nilpotent}: a product that lowers a bounded column
index more than $r$ times is zero.  Instantiate \texttt{exact\_hpl} and
\texttt{cycle\_transfer} for the double complex.
\end{enumerate}

For the first-order endpoint, encode an $m$ by $n$ matrix as its dimensions
and a finite list of signed numerator/positive-denominator pairs.  Define
dimension checks, entries, products, row sums, maximums and equality tests
as fixed primitive-recursive functions, with total default values on
malformed codes.  Guard every theorem by a Boolean validity test.  The
universally closed assertion is
\[
 \operatorname{ValidContraction}(w)=1\ \wedge\
 \operatorname{NormGuards}(w,q)=1
 \ \Longrightarrow\
 \operatorname{ResidualBounds}(w,N,q)=1.
\]
Bounded array assertions are represented by Boolean folds.  Their prefix
identities and norm inequalities use quantifier-free induction; the
denominator margin is explicit.  A proof of this sentence in
\texttt{Prf PRAAx}, followed by the repository's established bridge, is the
formal iPRA endpoint.  Ordinary Lean proofs of matrix identities, the
Banach-space interpretation, and exact-rational executions support distinct
claims.  This chapter supplies the complete mathematical argument and
construction; it does not claim that this full object derivation has already
been implemented.

\section{Lean definitions, checked results, and the iPRA export}\label{sec:formalization}

\subsection{The checked standalone kernel}

The accompanying project was compiled with Lean 4.32.0. It imports only
the distribution's \texttt{Init} and \texttt{Std} libraries; no Mathlib
download or repository checkout is required to replay it. There are
58 named proved theorems and two additional constructors whose numerical
contract fields are proved. The axiom audit prints 60 reports. Its actual
axiom union is \texttt{propext}, \texttt{Classical.choice}, and
\texttt{Quot.sound}. There are no admitted proofs or added correctness
axioms in the checked project. These facts are recorded in
\texttt{validation/lean\_status.json} and the full replay log.

\begin{longtable}{@{}p{0.25\textwidth}p{0.08\textwidth}p{0.60\textwidth}@{}}
\toprule
Module & Proofs & Exact checked scope\\
\midrule
\endhead
\texttt{FiniteMeasure} & 24 & Rational sums with internally variable finite lengths; finite Fubini; positivity; weighted errors; simultaneous perturbation of sample values and probability masses.\\
\texttt{Truncation} & 16 & Clipping bounds and nonexpansiveness, excess-tail formula and stability, scalar and finite weighted Vitali inequalities, and a three-term precision budget.\\
\texttt{Presentations} & 10 & Array-based finite probability data, refinement pullbacks, finite $L^1$ laws, and invariance of integrals, $L^1$ distance, capped distance, and tail integrals.\\
\texttt{ActionCompletion} & 8 & A concrete gain-dependent sampling schedule, regularity of mapped names, preservation of the same-name condition, and the proved fields of action composition.\\
\texttt{VitaliBlueprint} & 0 & Typechecked numerical definitions and target propositions for the stronger limit-producing theorem; the target propositions are not asserted as theorems.\\
\bottomrule
\end{longtable}

The checked namespace is \texttt{IPRAContinuousMeasure}. Representative
entry points are:
\begin{itemize}
\item expectation errors: \path{expectation_error}, \path{mass_perturbation};
\item clipping estimates: \path{finite_vitali_unit}, \path{vitali_budget};
\item refinement: \path{refinement_integral}, \path{refinement_l1},\\
      \path{refinement_rho}, \path{refinement_tail};
\item name completion: \path{action_regular}, \path{mapName_respects}.
\end{itemize}
The definitions \texttt{mapName} and
\texttt{composeDenseAction} contain checked proofs of their required
stability fields. The latter result is composition of metric actions;
linearity, positivity, the full stochastic tensor construction, and the
analytic theorems remain the report's mathematical results and development
targets rather than checked claims of this standalone core.

The project also kernel-checks four exact examples, including the
counterexample to \eqref{eq:false-old-vitali}. The separate Python programs
exercise larger finite cases. Their assertions are implementation checks;
the Lean universal statements and the mathematical proofs have separate
roles.

\subsection{All presentation data are explicit}

The interfaces are actual numerical structures rather than opaque
parameters named ``correctness.'' The following summary identifies their
fields and mathematical definitions; the full Lean definitions are included.

\begin{description}[style=nextline,leftmargin=0pt,labelwidth=0pt,labelsep=0pt]
\item[\texttt{FiniteProbability}]
An \texttt{Array Rat} of weights, a proof of nonnegativity of its in-range
entries, and a proof that its explicit rational fold sums to one. Array
access uses a total default outside the stored length. The functions
\texttt{integral}, \texttt{l1}, and \texttt{rho} are the finite sums in
\eqref{eq:m-simple} and \eqref{eq:m-stable-observers}.

\item[\texttt{Refinement}]
An \texttt{Array Nat} giving a parent for every fine atom, its range
proof, and the equation that the sum of each parent's descendant masses
equals the coarse mass. \texttt{Refinement.pull} is literal composition
of a value table with that parent array. Zero masses are allowed.

\item[\texttt{ProbabilityTower} and \texttt{SimpleTable}]
A numerical family of finite probabilities with coherent refinements;
a simple table is an array whose length matches the probability at its
stored stage. \texttt{SimpleTable.distance} lifts two arrays to the
maximum of their levels and evaluates the weighted absolute difference.
\texttt{measureDistance} and \texttt{tailIntegral} are similarly explicit.

\item[\texttt{RegularL1Name} and \texttt{SameL1Name}]
A sampler into simple tables and the bound
$d_1(s_k,s_\ell)\le2^{-k}+2^{-\ell}$.
The same-name predicate uses the corresponding cross-comparison at
arbitrary precisions. A full quotient/setoid API can be developed from
these data; it has not been silently imported as a new arithmetic sort.

\item[\texttt{DenseAction}]
A routine returning a simple output for a simple input and a precision,
a nonnegative rational gain, and the finite stability inequality.
An integer majorant for the gain is supplied separately to the name
construction and composition routines. \texttt{actionSample} uses a
precision offset based on that majorant.
The supplied implementation chooses a coarse integer offset; a proved
binary-length offset can improve its cost without changing the interface.

\item[\texttt{VitaliData}]
The uniformly named family, \texttt{UniformTailControl}, and
\texttt{InMeasureCauchyRate}, including the monotonicity and numerical
error contracts. \texttt{ceilingLog2} is a bounded search, not a real
logarithm. \texttt{vitaliPrecision}, \texttt{vitaliIndex}, and
\texttt{vitaliSample} are the actual programs in
\eqref{eq:m-vitali-rate}--\eqref{eq:m-output-name}.
\end{description}

\texttt{VitaliNameRegularityGoal}, \texttt{VitaliLiftingGoal}, and
\texttt{VitaliInheritedTailGoal} state the remaining propositions for the
constructed samples. A definition of type \texttt{Prop} records a target;
it supplies no proof of it. The analytic theorem does not assume an
integrable limit, and neither does \texttt{VitaliLiftingGoal}.

The host Lean structures permit arbitrary Lean functions as fields.
An iPRA instance fixes PR implementations for those fields and supplies
their object-theory contracts. Keeping this step explicit prevents a
higher-type semantic interface from being misreported as a uniform
first-order theorem about all arbitrary functions.

\subsection{An explicit object-language Vitali endpoint}

Here is a direct denominator-cleared target, sufficient to remove any
ambiguity about the proposed arithmetic translation. Let $w_i$ be natural
weights. Represent real samples $x_i,y_i$ with one positive denominator
$D$ and signed integer numerators $X_i,Y_i$. Let the integer cutoff
$\kappa\ge D$ represent $K=\kappa/D\ge1$. Define natural quantities
\[
 d_i=|X_i-Y_i|,\quad
 u_i=(|X_i|-\kappa)^+,\quad v_i=(|Y_i|-\kappa)^+,\quad
 c_i=\min(D,d_i).
\]
All can be written using natural truncated subtraction. If
$X_i=X_i^+-X_i^-$ and $Y_i=Y_i^+-Y_i^-$, for example,
\[
 d_i=((X_i^++Y_i^-)\mathbin{\dot-}(X_i^-+Y_i^+))
      +((X_i^-+Y_i^+)\mathbin{\dot-}(X_i^++Y_i^-)),
\]
where $a\mathbin{\dot-}b=\max(a-b,0)$; the parentheses indicate the two
complete minuends. Define the other absolute values in the same way.
The finite theorem to export is
\begin{equation}\label{eq:object-vitali}
 D\sum_{i<n}w_id_i
 \le D\sum_{i<n}w_i(u_i+v_i)
       +2\kappa\sum_{i<n}w_ic_i.
\end{equation}
Dividing by $D^2\sum_iw_i$ when that total is positive is precisely the
normalized clipping inequality. The unnormalized integer inequality
itself also holds when every weight is zero.

\paragraph{Definitions and proof program.}
Encode the four signed-numerator arrays, the weights, $n,D,\kappa$ in a
finite tuple. Define each lookup with a length guard and a zero default.
Define the three sums by primitive recursion on their prefix length $j$.
The invariant is the displayed inequality with $n$ replaced by $j$.
It is quantifier-free. At $j=0$ it is $0\le0$. At $j+1$ add the scalar
inequality
\[
 Dw_jd_j\le Dw_j(u_j+v_j)+2\kappa w_jc_j,
\]
which follows from the same finite clipping cases as
Lemma~\ref{lem:m-clipping}. Every branch compares natural integers.
The code-validity and cutoff guards remain explicit assumptions.
The theorem is then universally closed over these numerical parameters.

\paragraph{Reuse in the existing backend.}
The repository supplies the needed target syntax and rule
\texttt{PRA.qfInd} in \texttt{ConciseIPRA.PR}; natural sum/order and
signed-integer arithmetic in \texttt{IPRADerivations.Core}; rational
normalization, finite-sum semantics and monotonicity in
\texttt{Analysis.RationalFiniteSum}; and a closely related prefix
truncation proof in \texttt{Probability.TruncationEstimate}
\cite{repo-pr,repo-rational-sum,repo-truncation}.
The new file should construct a proof value of the form
\[
 [\,]\ \vdash_{\mathtt{PRAAx}}\
 \forall\vec a\,\bigl(\operatorname{Valid}(\vec a)=1
       \longrightarrow\operatorname{VitaliClear}(\vec a)=1\bigr),
\]
where \texttt{VitaliClear} is the characteristic function of
\eqref{eq:object-vitali}. This is a real proof-construction task. The
standalone theorem over \texttt{Rat} is a checked mathematical guide to it,
not an automatic quotation into \texttt{Prf PRAAx}.

\paragraph{From the finite endpoint to the named theorem.}
For fixed PR routines $S,U,R$, the desired object sentence universally
quantifies numerical indices and precisions. Its premises are the actual
regularity, tail, and capped-distance contracts of those routines. Apply
\eqref{eq:object-vitali} to finite common-refinement samples; use
\eqref{eq:m-observer-lipschitz} to account for their errors; substitute the
bounded logarithm program; and prove regularity of
$b_k=S(N(k+2),k+2)$ by the explicit three-error calculation.
The final convergence statement is expressed using cross-distances of
the two samplers with arbitrary observation errors. It need not quantify
over an uncountable space of limits. The inherited-tail property follows
by an arbitrary-precision comparison, not by a full induction schema for
analytic formulas.

\subsection{Development order for the remaining mathematics}

\begin{longtable}{@{}p{0.22\textwidth}p{0.39\textwidth}p{0.32\textwidth}@{}}
\toprule
Target & Necessary definitions and induction data & Route to completion\\
\midrule
\endhead
Full Vitali lifting & Common-refinement pseudometric laws; regular names; tail and capped-distance stability; bounded ceiling logarithm & Prove the typed goals in \texttt{VitaliBlueprint}; then implement \eqref{eq:object-vitali} and its clock substitutions in the object backend.\\
Conditional expectation & Cell numerators, exact zero-weight branch, parent-map composition, integral-characterization equation & Reuse \texttt{TowerCommonDen.regroup}; prove contraction after canceling cell mass, then lift and prove \eqref{eq:m-upward-rate}.\\
Stochastic operator category & Linear/positive/unital/integral action fields; strong-action sampler; atom-product tables & Extend the checked metric-action constructors with the stated algebraic fields. Prove equality on dense inputs first; tensor by the indicator estimate.\\
Integrated formulas & Many-sorted finite syntax with bound-variable indices; interval and Lipschitz annotations; net and quadrature codes & Structural recursion for the evaluator and budgets; induction on the fixed syntax; array induction for each numerical binder. Clip intermediate values to their annotated ranges.\\
Holomorphic integrals & Complex rational pairs; coefficient-name family; radius/majorant data; finite derivative sums & Prove finite geometric and polynomial remainder identities, then use their explicit tail clocks and the integral contraction.\\
Homological transfer & Graded matrix records, normalized contraction equations, perturbation equation, row-sum norm, finite Neumann sums & Follow the exact identity and residual route in Section~\ref{sec:homological}; use its matrix blueprint and denominator-cleared bounded folds.\\
\bottomrule
\end{longtable}

For an internally variable syntax size or matrix dimension, the compiler
must use numerical arrays and bounded characteristic folds, rather than
generate one externally unrolled theorem per size. The repository's current
bounded-syntax and rational-array providers are suitable starting points.
For a fixed mathematical formula, a schema indexed externally by that
formula is already meaningful; it should simply be labeled as such.

\subsection{Reproduction and verification limits}

Run \texttt{python3 verify.py} for the exact finite demonstrations. With
the pinned Lean toolchain installed, run \texttt{bash lean/replay.sh} from
the appropriate working directory, or use the supplied Makefile target.
\texttt{python3 build\_pdf.py} rebuilds this document. The TeX helper
can create local formats and font maps from installed TeX Live resources
on a minimally configured system; it does not download missing packages.

The manifest records the source commit and Git blob identifiers of the
upstream files actually used. Those files were inspected, not rebuilt in
this session. The standalone Lean replay, exact Python checks, and document
build were performed here. External toolchain binaries and build caches
are excluded from the archive. No claim of a complete iPRA formalization
of the whole analytic framework follows from the 58 host theorems.

\section{Resolved difficulties and the remaining mathematical boundary}\label{sec:resolved}

The framework provides more than a uniform instruction to supply moduli.
Several of the required data are \emph{derived} from existing presentations:
a dominating $L^1$ name supplies a tail cutoff; a regular simple name
supplies a conditional-expectation convergence rate; a compact net and
distance-function names supply continuous tent quadrature; a contraction
and a perturbation gap supply inverse and homotopy clocks. Each derivation
has an explicit finite routine and a proof of its numerical contract.

\begin{longtable}{@{}p{0.28\textwidth}p{0.40\textwidth}p{0.25\textwidth}@{}}
\toprule
Difficulty & Resolution obtained & Remaining hypothesis\\
\midrule
\endhead
Incorrect separate-tail truncation in the earlier Vitali proof & Replace it by the checked clipping estimate, then build the limit from pairwise Cauchy data. & A uniform cutoff program and a capped-distance Cauchy rate.\\
The limit is assumed integrable before convergence is proved & The diagonal simple-code sampler constructs the $L^1$ limit and derives its uniform tails. & Joint PR names for the sequence.\\
An additional amplitude-tail oracle appears necessary under domination & Extract a cutoff from the finite supremum of each simple approximation to the dominator. & A regular $L^1$ name for the dominator.\\
Small cells cause large posterior values and unstable division & Use the full $L^1$ conditional operator; the cell mass cancels the reciprocal. & A positive mass margin for a pointwise posterior query.\\
Finite operators fail to approximate the identity uniformly & Compose strong input-dependent actions with certified gain schedules. & An action routine on every supplied dense input; no automatic operator-name extraction.\\
Integration is absent from compact model evaluation & Add rational tent quadrature and an expectation binder with weighted errors. & Compact nets, compatible measure data, Lipschitz/range bounds.\\
An analytic integral has a singular spatial factor & Use $L^1$ coefficient bounds and integrate the norm-convergent series. & A radius gap and a jointly presented coefficient family.\\
Ordinary nerve methods require acyclic overlaps & Retain local chain models and transfer their full differential by finite homological perturbation. & Explicit local models and a valid cover-to-total-complex comparison.\\
An approximate differential is mistakenly treated as square-zero & Carry the exact residual identities and a geometric bound until the completed differential or exact nilpotent formula is available. & Exact square-zero of the original perturbed differential.\\
\bottomrule
\end{longtable}

These hypotheses cannot all be erased. The bounded martingale construction
in Section~\ref{sec:measure} makes a uniform computable convergence-rate
extractor decide halting. Small $L^1$ error cannot control a supremum.
Independent fair events prevent a compact-net treatment of an entire
atomless probability algebra. Small differential perturbations can change
homology; the Schur-complement example shows the change on an explicit
retained model. Exact real equality, unbounded disintegration, arbitrary
model selection, and unrestricted countable unions remain outside the
proved construction unless their missing data are supplied or derived.

For further repository work, the highest-value next object-language step
is the integer clipping inequality \eqref{eq:object-vitali}, followed by
the regular-name diagonal and conditional-average contraction. They use
the existing arithmetic backend and produce reusable continuous mathematics
with a clear verification boundary. The matrix residual endpoint is a
parallel target with similarly explicit bounded-sum invariants. Finishing
these endpoints would turn the present full mathematical arguments and
checked finite kernels into new, inspectable \texttt{Prf PRAAx} exports.

\begin{thebibliography}{99}
\small
\bibitem{neri-pischke}
Morenikeji Neri and Nicholas Pischke. \emph{Proof mining and probability theory}. Forum of Mathematics, Sigma 13 (2025), e187. arXiv:2403.00659v2.
\url{https://arxiv.org/abs/2403.00659}.

\bibitem{avigad-dean-rute}
Jeremy Avigad, Edward T. Dean, and Jason Rute. \emph{A metastable dominated convergence theorem}. Journal of Logic \& Analysis 4:3 (2012), 1--19. Preprint title: \emph{Metastable convergence theorems}.
\url{https://arxiv.org/abs/1108.4400}.

\bibitem{fritz}
Tobias Fritz. \emph{A synthetic approach to Markov kernels, conditional independence and theorems on sufficient statistics}. Advances in Mathematics 370 (2020), 107239.
\url{https://arxiv.org/abs/1908.07021}.

\bibitem{berenstein-henson}
Alexander Berenstein and C. Ward Henson. \emph{Model theory of probability spaces}. arXiv:2302.01519 (2023).
\url{https://arxiv.org/abs/2302.01519}.

\bibitem{bagheri-pourmahdian}
Seyed-Mohammad Bagheri and Massoud Pourmahdian. \emph{Continuous integration logic}. arXiv:1910.00191 (2019).
\url{https://arxiv.org/abs/1910.00191}.

\bibitem{crainic2004}
Marius Crainic. \emph{On the perturbation lemma, and deformations}. arXiv:math/0403266 (2004).
\url{https://arxiv.org/abs/math/0403266}.

\bibitem{repo-overview}
Vladimir Reshetnikov et al. Primitive recursive arithmetic: overview.
\repo, pinned commit \texttt{b2341af5370c}. Source inspected 20 November 2026.
\url{https://github.com/VladimirReshetnikov/ProveIt/blob/b2341af5370cd8a43f72ebb2341c0d7d7395b634/Logic/PrimitiveRecursiveArithmetic/README.md}.

\bibitem{repo-derivations}
Vladimir Reshetnikov et al. Object-theory derivations: module inventory.
\repo, pinned commit \texttt{b2341af5370c}. Source inspected 20 November 2026.
\url{https://github.com/VladimirReshetnikov/ProveIt/blob/b2341af5370cd8a43f72ebb2341c0d7d7395b634/Logic/PrimitiveRecursiveArithmetic/Derivations/README.md}.

\bibitem{repo-pr}
Vladimir Reshetnikov et al. The ordinary iPRA syntax and axioms.
\repo, pinned commit \texttt{b2341af5370c}. Source inspected 20 November 2026.
\url{https://github.com/VladimirReshetnikov/ProveIt/blob/b2341af5370cd8a43f72ebb2341c0d7d7395b634/Logic/PrimitiveRecursiveArithmetic/ConciseIPRA/Lean/ConciseIPRA/PR.lean}.

\bibitem{repo-abstract}
Vladimir Reshetnikov et al. Abstract and continuous mathematics: package description.
\repo, pinned commit \texttt{b2341af5370c}. Source inspected 20 November 2026.
\url{https://github.com/VladimirReshetnikov/ProveIt/blob/b2341af5370cd8a43f72ebb2341c0d7d7395b634/Logic/PrimitiveRecursiveArithmetic/docs/proofs/foundations/ipra_abstract_analysis_2026-11-20/README.md}.

\bibitem{repo-continuous-two}
Vladimir Reshetnikov et al. From finite observations to continuous structures: package description.
\repo, pinned commit \texttt{b2341af5370c}. Source inspected 20 November 2026.
\url{https://github.com/VladimirReshetnikov/ProveIt/blob/b2341af5370cd8a43f72ebb2341c0d7d7395b634/Logic/PrimitiveRecursiveArithmetic/docs/proofs/foundations/ipra_continuous_framework_2026-11-20_2/README.md}.

\bibitem{repo-complex-series}
Vladimir Reshetnikov et al. Earlier quantitative complex-series interface.
\repo, pinned commit \texttt{b2341af5370c}. Source inspected 20 November 2026.
\url{https://github.com/VladimirReshetnikov/ProveIt/blob/b2341af5370cd8a43f72ebb2341c0d7d7395b634/Logic/PrimitiveRecursiveArithmetic/docs/proofs/foundations/ipra_abstract_analysis_2026-11-20/sections/complex.tex}.

\bibitem{repo-expectation-packet}
Vladimir Reshetnikov et al. Expectation packet: convergence chapter.
\repo, pinned commit \texttt{b2341af5370c}. Source inspected 20 November 2026.
\url{https://github.com/VladimirReshetnikov/ProveIt/blob/b2341af5370cd8a43f72ebb2341c0d7d7395b634/Logic/PrimitiveRecursiveArithmetic/docs/proofs/probability/expected_value_ipra_2026-11-16/paper/sections/12_convergence.tex}.

\bibitem{repo-finite-kernel}
Vladimir Reshetnikov et al. Finite expectation coding.
\repo, pinned commit \texttt{b2341af5370c}. Source inspected 20 November 2026.
\url{https://github.com/VladimirReshetnikov/ProveIt/blob/b2341af5370cd8a43f72ebb2341c0d7d7395b634/Logic/PrimitiveRecursiveArithmetic/Derivations/Lean/IPRADerivations/Probability/FiniteKernel.lean}.

\bibitem{repo-tower-common-den}
Vladimir Reshetnikov et al. Common denominators and the finite tower law.
\repo, pinned commit \texttt{b2341af5370c}. Source inspected 20 November 2026.
\url{https://github.com/VladimirReshetnikov/ProveIt/blob/b2341af5370cd8a43f72ebb2341c0d7d7395b634/Logic/PrimitiveRecursiveArithmetic/Derivations/Lean/IPRADerivations/Probability/TowerCommonDen.lean}.

\bibitem{repo-truncation}
Vladimir Reshetnikov et al. Common-dominator truncation estimate.
\repo, pinned commit \texttt{b2341af5370c}. Source inspected 20 November 2026.
\url{https://github.com/VladimirReshetnikov/ProveIt/blob/b2341af5370cd8a43f72ebb2341c0d7d7395b634/Logic/PrimitiveRecursiveArithmetic/Derivations/Lean/IPRADerivations/Probability/TruncationEstimate.lean}.

\bibitem{repo-fubini}
Vladimir Reshetnikov et al. Finite Fubini in object iPRA.
\repo, pinned commit \texttt{b2341af5370c}. Source inspected 20 November 2026.
\url{https://github.com/VladimirReshetnikov/ProveIt/blob/b2341af5370cd8a43f72ebb2341c0d7d7395b634/Logic/PrimitiveRecursiveArithmetic/Derivations/Lean/IPRADerivations/Core/Fubini.lean}.

\bibitem{repo-rational-sum}
Vladimir Reshetnikov et al. Numerical rational finite sums.
\repo, pinned commit \texttt{b2341af5370c}. Source inspected 20 November 2026.
\url{https://github.com/VladimirReshetnikov/ProveIt/blob/b2341af5370cd8a43f72ebb2341c0d7d7395b634/Logic/PrimitiveRecursiveArithmetic/Derivations/Lean/IPRADerivations/Analysis/RationalFiniteSum.lean}.

\bibitem{repo-dag}
Vladimir Reshetnikov et al. Rational error budgets on finite dependency graphs.
\repo, pinned commit \texttt{b2341af5370c}. Source inspected 20 November 2026.
\url{https://github.com/VladimirReshetnikov/ProveIt/blob/b2341af5370cd8a43f72ebb2341c0d7d7395b634/Logic/PrimitiveRecursiveArithmetic/Derivations/Lean/IPRADerivations/Foundations/PrecisionCostDAG.lean}.

\bibitem{repo-dag-schedule}
Vladimir Reshetnikov et al. An internal precision schedule for dependency graphs.
\repo, pinned commit \texttt{b2341af5370c}. Source inspected 20 November 2026.
\url{https://github.com/VladimirReshetnikov/ProveIt/blob/b2341af5370cd8a43f72ebb2341c0d7d7395b634/Logic/PrimitiveRecursiveArithmetic/Derivations/Lean/IPRADerivations/Foundations/PrecisionCostDAGSchedule.lean}.

\bibitem{repo_hpl}
Vladimir Reshetnikov et al. Unknot recognition, report 13: finite perturbation theory.
\repo, pinned commit \texttt{b2341af5370c}. Source inspected 20 November 2026.
\url{https://github.com/VladimirReshetnikov/ProveIt/blob/b2341af5370cd8a43f72ebb2341c0d7d7395b634/Topology/UnknotRecognition/reports/13/paper/perturbation_theory.tex}.

\end{thebibliography}

\end{document}
